\input{preamble}



\begin{document}



\title{Seminars in Differential Geometry}

\maketitle



\phantomsection

\label{section-phantom}



\tableofcontents



\section{A primer on symplectic groupoids}
\label{section-prime-sympl-group}

\noindent
Camilo Angulo, Jilin University.
Geometric Structures Seminar, IMPA.
February 13, 2025.

\medskip\noindent
{\bf Abstract.}
In the late 17th century, Simeon Denis
Poisson discovered an operation that
helped encoding and producing conserved
quantities. This operation is what we
now know as a Lie bracket, an infinitesimal
symmetry, but what is its global
counterpart? Symplectic groupoids are one
possible answer to this question. In
this talk, we will introduce all the basic
concepts to define symplectic
groupoids, and their role in Poisson
geometry. We will discuss key examples, and
applications. The talk will be accessible to
those familiar with differential
geometry, but no prior knowledge of groupoids
will be assumed.


\medskip\noindent
{\bf Part 1. Poisson geometry.}

Hamiltonian formalism. Recall that being a
conserved quantity \(f \in
C^\infty(X)\) is the same thing as
\(\{H,f\}=0\).

\begin{itemize} \item We have seen that it is
always possible to take quotient
of a symplectic manifold with a group action
to obtain a  Poisson
manifold.  \item Then we have found a way to
produce a symplectic foliation
from a 2-vector \(\pi \in
\mathfrak{X}^2(M):=\Lambda^{2}(TM)\).  \item

\begin{remark}
\[\left\{ \text{Lie algebra on
\(\mathfrak{g}\)}
\right\} \xrightarrow{1-1}\left\{
\text{Linear Poisson bracket on
\(\mathfrak{g}^*\)}  \right\}  \]
\end{remark}

\item We saw very nice examples
of foliation that have to do with Lie
algebras. So \(\mathfrak{b}^*_3\) which
gives the ``open book foliation", and
\(\mathfrak{e}^*\) that gives a foliation
by cylinders.
\end{itemize}

\medskip\noindent
{\bf Part 2. Symplectic realizations.}

Consider
\[(\Sigma,\omega) \xrightarrow{\mu}(M,\pi)\]


\iffalse\begin{tikzcd}
	T_p\Sigma
	\arrow[r,"d_p\mu"]\arrow[d,"\omega^\flat"]&
	T_xM\\
	T_p^*\Sigma &
	T^*_pM\arrow[l,"(d\mu)^*",swap]\arrow[u,"\pi
	^\sharp"]
\end{tikzcd}\]\fi
So that
\[\pi ^\sharp=d_p\mu \circ \omega^{-1}
\circ(d\mu)^*\]



\begin{lemma}
\(\dim (\Sigma) \geq  2 \dim (M) -
\operatorname{rk}(\pi_x)\) for all \(x \in
M\).
\end{lemma}

\begin{proof}
Done in seminar.
\end{proof}

\begin{example}
\((\mathbb{R}^2,0)\). So the map
$$
\begin{aligned}
	(\mathbb{R}^4,dx \wedge du + dy \wedge dv)
	&\longrightarrow  \mathbb{R}^2\\
	(x,y,u,v) &\longmapsto (x,y)
\end{aligned}
$$

\end{example}

\begin{exercise}
Find the symplectic realization \(\omega\) in
\((\mathbb{R}^4, \omega )
\xrightarrow{\mu}(\mathbb{R}^2,(x^2+y^2)
\partial_x
\wedge \partial_y\)

$$
\begin{aligned}
	(\mathbb{R}^4, \omega) &\longrightarrow
	(\mathbb{R}^2,(x^2+y^2) \partial_x \wedge
	\partial _y) \\
	(x,y,z,w) &\longmapsto (x,y)
\end{aligned}
$$

Also find the symplectic realization of
\(\operatorname{ a f f }^*\) with
\(\{x,y\}=x\).
\end{exercise}

\medskip\noindent
{\bf Part 3: Groupoids}

\medskip\noindent
{\bf Motivation}
\begin{enumerate}
\item Fundamental grupoid: objects are points
in the manifold and arrows are paths.
\item The action of \(S^1\) on \(S^2\) by
rotation does not give a nice quotient
because there are two singular points.
Consider the groupoid  \(S^1 \times S^2\) of
orbits. These are the arrows. The points are
just the points of \(S^2\).
\item Consider a foliation (like Möbius
foliation of circles; where there is a
singular circle, the soul). You can do the
same thing as in fundamental groupid
leafwise. Arrows then are equivalence classes
of paths that live inside leaves. This is
called monodromy of a foliation. Again
objects are points.
\item You can take a quotient of monodromy
using a connection given by the foliation.
This allows to identify certain paths between
the leaves. This is called {\it holonomy (of
a foliatiation)}. (So you can make this
notion match the usual holonomy given by
riemmanian connection.)

{\bf Upshot.}
	So the point is taking some sort of function
	space on these groupoids you can gather the
	information given by the non-smooth quotient
	(like in the case of the sphere rotating).
	So this groupoid motivation says how to get
	some structure that resembles a non smooth
	quotient.
\item Last motivation: the grid of squares
has a tone of symmetries. If you restrict to
just a few squares you loose so many
symmetries. But there's a grupoid hidden in
there that tells you what you intuition knows
about this finite grid of squares.
\end{enumerate}

\begin{definition}
A {\it groupoid} is a category where all
morphisms are invertible.
\end{definition}

So there is a kind of product among the
objects, given by composition but:
 not every two pair of objects can be
 multiplied!---only those whose
source and target match. So that's the lance
about grupoids.

Just so you make sure you understand: the
groupoid \(G\) is the morphisms of the
category. The objects are points (of a
manifold).

\begin{definition}
{\it Lie grupoid} is when the following
diagram is inside category of smooth
manifolds and \(s,t\) (source and target
maps) are submersions:
$$
\xymatrix{
G^{(2)}\ar[r]^m& G\ar[r]^ s_t & M\ar[r]^u&G
}
$$
\end{definition}

\begin{proposition}[Properties of Lie
groupoids]
\begin{itemize}
\item \(m\) is also a submersion.
\item \(i\) (inversion) is a diffeomorphism.
\item  \(u\) (unit=identity) is an embbeding.
\end{itemize}
\end{proposition}

\begin{definition}
Consider \(x \in M\) and the inverse image of
source map: \(s^{-1}(x)=\{\text{arrows that
start at \(x\)} \}\). Now if you act with
\(t\) on this set you get {\it the orbit} of
\(x\): \(\{\text{\(y \in M\) such that there
is an arrow from \(x\) to \(y\)} \}\).

And there also {\it an isotropy} \(G_x=\{g
\in G: \text{\(g\) goes from \(x\) to \(x\)}
\}\)
\end{definition}

\begin{example}
\begin{enumerate}
\item \(G=M\),  \(M=M\).
 \item Lie groups.
	\item Lie group bundles.
\item \(G=M \times M\), \(M =M\).
 \item Fundamental groupoid. Isotropy group
 is fundamental group! And orbit
 is…\clearpage universal cover!
\item Subgroupoids.
\item  Foliations.
\item If you have a normal group action of
\(G\) on \(M\) you construct a groupoid
action with groupoid \(G \times M\) and
objects \(M\), with product given on the
group part of the product. Orbits are orbits.
Isotropy group is isotropy group.
 \item Principal bundles.
\end{enumerate}
\end{example}

{\bf Back to Poisson.}

There's also a notion of Lie algebroid. Which
is strange. But the point is that to every
Poisson manifold there is a Lie algebroid.

So the question is whether there is a Lie
groupoid associated to that Lie algebroid.
Not always.

\medskip\noindent
{\bf Big question}[Fernandez and ?]
When a symplectic manifold is integrable?


(Remember that integrating means go from
algebra(oid) to group(oid).

And the point is that:

\medskip\noindent
When you {\it can} go back, you get a
\textit{ symplectic groupoid}.

\begin{remark}
Look for Kontsevich's notes on Weinstein!
\end{remark}

\begin{remark}
History: Weinstein did this intending to do
quantization (geometric?) on Poisson
manifolds. (That involves a \(C^*\)algebra
coming from the symplectic groupoid.)
\end{remark}

\begin{definition}
A {\it symplectic groupoid} is a groupoid
\(G, M\) together with  \(\omega \in
\Omega^2(M)\) such that \(\omega\) is
symplectic and multiplicative, meaning that
\(\partial  \omega =0\), that is, \(\iff
m^*\omega=\operatorname{pr}_1^*\omega+
\operatorname{pr}^*_2\omega\in
\Omega^{2}(G^{(2)})\iff\) take two vectors \(
X_k,Y_k \in TG\), and \(\omega(X_0 \star
Y_0,X_1\star Y_1) =  \omega(X_0,Y_0) + \omega
(X_1,Y_1)\)
\end{definition}

\begin{theorem}
If \((G, \omega)\) is a symplectic groupoid,
then
\begin{enumerate}
\item \(\exists !\) poisson structure on
\(M\)
\item for which \(t: G \to M\) is a
symplectic realization,
\item  Leaves are connected components of
orbits,
\item \(\mathsf{Lie}(G)\cong T^*M\) via \(X
\mapsto -u ^*(i_X \omega\).
\end{enumerate}
\end{theorem}

\begin{remark}
Look for Alejandro Cabrera, Kontsevich. There
are two things one is de Rham and the other…
from the future: simplicial?
\end{remark}

\medskip\noindent
{\bf Upshot.} The obstruction to knowing when
symplectic groupoid exists is ``variation of
symplectic form \(\omega=(1+t^2)
\omega_{S^2}\)". So how does the symplectic
group vary from leaf to leaf. So there are
two situations in which the thing doesn't
work.

\section{Spheres with minimal equators}
\label{section-spheres-with-minimal-equators}

\noindent
Lucas Ambrozio, IMPA.
Differential Geometry Seminar, IMPA.
June 24, 2025.

\medskip\noindent
{\bf Abstract.} We will discuss the
connection between Riemannian metrics on the
sphere with respect to which all equators are
minimal hypersurfaces, and
algebraic curvature tensors with positive
sectional curvatures.

\bigskip\noindent

\begin{definition}
\label{definition-equator}
An {\it $(n-k)$-equator} orthogonal to $\Pi$
is
$$
\Sigma_\Pi:=\{p\in\mathbb{S}^n:\left\langle{}p,
x\right\rangle{}=0\forall
x\in\Pi\}
$$
for $\Pi$ a $k$-dimensional linear subspace
of $\mathbb{R}^{n+1}$.
\end{definition}

\begin{remark}
\label{remark-equat-total-geode-usual-metr}
Equators are totally geodesic hypersurfaces
with the usual sphere metric, which
implies they are minimal hypersurfaces.
\end{remark}

{\bf Problem.} Characterize the set
$\mathcal{M}_k(U)$ of metrics $g$ on an
open set $U\subset\mathbb{S}^n$ such that all
$k$-equators $\Sigma_\Pi$ with
 $\Sigma\cap U\neq\emptyset$ yield are
 minimal hypersurfaces $\Sigma\cap U$
on $(U,g)$.

\begin{remark}
\label{remark-why-not-Rn}
This problem can be thought of as a problem
of finding metrics on $\mathbb{R}^n$
such that $k$-planes are minimal. To see why
project the $k$-equators to
$T_p\mathbb{S}^n$ and pullback those metrics
to the sphere.
\end{remark}

\medskip\noindent

Let $g\in\mathcal{M}_k(U)$ for
$U\subset\mathbb{S}^n$ open and $n\geq 2$.

\begin{theorem}[Beltrami, Schäfli]
\label{theorem-Beltrami-Schafli}
If $k=1$ then $g$ has constant sectional
curvature.
\end{theorem}

\begin{theorem}[Hongan]
\label{theorem-Hongan}
If $1<k<n-1$ then $g$ has constant sectional
curvature.
\end{theorem}

Then Hongan also managed to produce a
classification of these metrics for
$k=n-1$.

\begin{remark}
\label{remark-linea-inver-prese-equat-mg}
If $T\in \text{GL}(n+1,\mathbb{R})$, then
$$
\begin{aligned}
\varphi: \mathbb{S}^n &\longrightarrow
\mathbb{S}^n \\
x &\longmapsto \frac{Tx}{|Tx|}
\end{aligned}
$$
is a diffeomorphism that maps $k$-equators
into $k$-equators. Thus if
$g\in\mathcal{M}_k(\mathbb{S}^n)$ then so is
$\varphi(T)^*g$.
\end{remark}

\begin{theorem}
\label{theorem-equivariant-bijection}
There exists a $\text{GL}(n+1,\mathbb{R})$
equivariant bijection
$$
\mathcal{M}_{n-1}(\mathbb{S}^n)
\leftrightarrow
\text{Curv}_+(\mathbb{R}^{n+1})
$$
where the set on the right-hand-side is the
set of algebraic curvature tensors
(also called curvature-like, i.e. with the
same symmetries as the Riemannian
curvature tensor) on $\mathbb{R}^{n+1}$ with
positive sectional curvature.

The group action is given as follows for
$T\in\text{GL}(n+1,\mathbb{R})$:
$$
(R\cdot
T)(x,y,z,w)=\frac{1}{|\det(T)|^{\frac{1}{n+1}
}}R(Tx,Ty,Tz,Tw)
$$
\end{theorem}

The point is that
$\text{Curv}_+(\mathbb{R}^{n+1})$ is an open
cone on a linear
space. Here are two simple corollaries:

\begin{lemma}
\label{lemma-corollaries}
\begin{enumerate}
\item $\mathcal{M}_{n+1}(\mathbb{S}^n)$ is in
bijection with an open positive
cone of an
$\frac{n(n+2)(n+1)^2}{12}$-dimensional real
vector space.
\item Every metric on
$\mathcal{M}_{n-1}(\mathbb{S}^n)$ is
invariant by the
antipodal map.
\end{enumerate}
\end{lemma}

{\bf Algorithm.} From any $R\in
\text{Curv}_p(\mathbb{R}^{n+1})$ we obtain a
symmetric positive definite
(positive-definitiness comes from the
positiveness
of the curvature of $R$) 2-tensor  $k_R$
satisfying
$$
(k_R)_p(v,v)=R(pv,pv)>0
$$
Also, $k_R$ has the  {\it Killing property},
i.e. that $\overline{\nabla}k(X,X,X)=0$ for
all
$X\in\mathfrak{X}(\mathbb{S}^n)$.

Then we define a positive function on
 $\mathbb{S}^n$ by
\begin{equation}
\label{equation-def-DR}
D_R:=\left(\frac{d\text{Vol}_{k_R}}{dV_g}
\right)^{\frac{4}{n-1}}
\end{equation}
and finally a Riemannian metric on
$\mathbb{S}^n$ in
$\mathcal{M}_{n-1}\mathbb{S}^n$ by
$$
g_R=\frac{1}{D_R}k_R
$$
And to go back, for $g \in
\mathcal{M}_{n-1}(\mathbb{S}^n)$ define a
positive
function on $\mathbb{S}^n$
$$
F_g:=\left(\frac{dV_g}{dV_{\overline{g}}}
\right)^{\frac{4}{n-1}}
$$
Then let $k_g:=\frac{1}{F_g}g>0$, which is a
positive definite Killing 2-tensor,
from which we may define
$R_g\in\text{Curv}_+(\mathbb{R}^{n+1})$ with
$R_g(pv,pv)=(k_g)_p(v,v)$ for all $p,v\in
T\mathbb{S}^n$.

More corollaries:
\begin{lemma}
\label{lemma-corollaries2}
\begin{enumerate}
\item $g\in\mathcal{M}_{n-1}(\mathbb{S}^n)$
is analytic because it is a Killing
tensor on $\mathbb{S}^n$, which are
well-known.
\item If $g$ is left-invatiant on
$\mathbb{S}^3$, seen as unit quaternions,
then
$g\in\mathcal{M}_2(\mathbb{S}^3$. Moreover,
for $a\geq b\geq c>0$,
$$
aL_i\odot L_i+bL_j\odot L_j+cL_k\odot L_k=k
$$
is Killing, $k>0$,  $D_k$ constant and thus
$g=\frac{1}{\text{const.}}k
\in\mathcal{M}_2(\mathbb{S}^3)$.
\item $R$ curvature tensor of
$(\mathbb{C}P^{2},g_{FS})$. We may not
remember
what's the curvature tensor, but we know the
sectional curvature is $1\leq
\text{sec}(R)\leq 4$,
$$
(k_R)_p(v,w)=\overline{g}(v,w)+3\overline{g}
(Jp,v)\overline{g}(Jp,w)
$$
and $D_R=4^{\frac{4}{3-1}}=4$, so that by
\ref{equation-def-DR} we obtain
$g_R=\frac{1}{4}k_R$, which is a Berger
metric on $\mathbb{S}^3$ with scalar
curvature 0.
\end{enumerate}
\end{lemma}

\bigskip
Now define
$$
\Sigma_V=\{p\in\mathbb{S}^n:\left\langle{}p,v\right\rangle{}=
0\}=V^{-1}(0)
$$
where $V(x):=\left\langle{}x,v\right\rangle{}$ for all
$x\in\mathbb{S}^n$. Then the normal
vector field is $\nabla V/|\nabla V|_g$, and
the second fundamental form is
given by
$$
A=\frac{1}{|\nabla V|_g}\text{Hess}_gV
$$
and its mean curvature by
\begin{equation}
\label{equation-mean-curva-level-sets-V}
H=\frac{1}{|\nabla
V|_g}\left(\Delta_gV-\text{Hess}
_gV\left(\frac{\nabla
V}{|\nabla V|},\frac{\nabla V}{|\nabla
V|}\right)\right)
\end{equation}
For every $v \in \mathbb{S}^n$ and
$p\in\Sigma_V$, we see that $H_{\Sigma_V}=0$
iff
$$
|\nabla
V|^2_g(p)\Delta_gV(p)-\text{Hess}_gV(\nabla
V(p),\nabla V(p))=0
$$
And for $\overline{g}$,
$$
\text{Hess}_{\overline{g}}V+V\overline{g}=
0\implies
\text{Hess}_{\overline{g}}V(X,X)=0
$$
for all $X\in T_p\mathbb{S}^n$ and
$p\in\Sigma_v$. Then
$$
J_g(X,Y,Z)=g(\nabla_XY-\overline{\nabla}_XY,
Z)
$$
$$
J_g(X,Y,\nabla
V)=\text{Hess}_{\overline{g}}-\text{Hess}
$$
{\bf Problems.}
\begin{enumerate}
\item Similar story for
$\mathbb{C}P^{n},\mathbb{H}P^n$?
\item Complete metrics on $\mathbb{R}^n$ with
minimal hyperplanes.
\item Find geometric invariants of metrics on
$\mathcal{M}_{n-1}(\mathbb{S}^n)$
(may be useful to study $(M^n,g)$, $n\geq 4$,
$\text{sec}>0$.
\end{enumerate}

\section{On wrapped Floer homology barcode
entropy and hyperbolic sets}
\label{section-wrapp-floer-homol-barco}

\noindent
Rafael Fernandes, UC Santa Cruz.
Differential Geometry Seminar, IMPA.
September 4, 2025.

\medskip\noindent
{\bf Abstract.} In this talk, we will discuss
the interplay between the
wrapped Floer homology barcode and
topological entropy. The concept of barcode
entropy was introduced by Çineli, Ginzburg,
and Gürel and has been shown to be
related to the topological entropy of the
underlying dynamical system in various
settings. Specifically, we will explore how,
in the presence of a topologically
transitive, locally maximal hyperbolic set
for the Reeb flow on the boundary of
a Liouville domain, barcode entropy is
bounded below by the topological entropy
restricted to the hyperbolic set.

\medskip\noindent

Let $M^n$ be a manifold.
$\omega \in \Omega^2(M)$ is {\it symplectic}
if $d \omega=0$ and it is
nondegenerate.

\begin{example}
\label{example-Rn-is-symplectic}
$\mathbb{R}^n$ is symplectic with canonical
Darboux form.
\end{example}

Recall the definition of Hamiltonian vector
field
associated to a function $H \in C^\infty(M)$.

\begin{definition}
\label{definition-nonde-diffe}
A diffeomorphism $\varphi:M \to M$ is called
{\it non-degenerate} if
$\Phi(\varphi)\cap\Delta \subset M\times M$
(pitchfork, i.e. transversal intersection!).
\end{definition}

Let $M^{2n}$ be a closed symplectic manifold.
Arnold's conjecture says
\begin{enumerate}
\item If $\varphi=\varphi_H$ (Hamiltonian
flow) is nondegenerate, then
$$
\# \text{Fix}(\varphi_H) \geq
\sum_{i=0}^{2n}\dim H_i(M,k)=\dim H_*(M,k)
$$
\item If $\varphi=\varphi_H$ is degenerate,
then
$$
\# \text{Fix}(\varphi)\geq \text{Cl}(M)+1
$$
where $\text{Cl}(M)$ is the maximum number of
homology classes
we can add before getting to zero.
\end{enumerate}

Why do we care? Because
$$
\#
\text{Fix}(\varphi_H)
\leftrightarrow\{\text{1-periodic
orbits of $X_H$}\}
$$

Idea by Floer. Construct an invariant that
would
say something about periodic orbits.

\medskip\noindent
{\bf Question.} Can Floer theory capture
other ``dynamical information''?
(Other than the periodic orbits.)

\medskip\noindent
A {\it persistence module} is a pair
$(V,\Pi)$, where $V=\{V_t\}_{t \in
\mathbb{R}}$ is a family of
$\mathbb{F}$-vector spaces
and $\Pi=\{\Pi_{st}\}_{s \leq  t}$ is a
family of maps such that
\begin{enumerate}
\item $\Pi_{ss}=, \Pi_{ts}\circ
\Pi_{rt}=\Pi_{rs}$.
\item $\exists  s \subset \mathbb{R}$ such
that
$\Pi_{s t}$ is an isomorphism for $s,t$ in
the same connected component of
$R\setminus S$.
\item $\Pi_{s t}$ have finite rank.
\item $\exists s_0$ $V_s=\{0\}$, $s \leq
s_0$.
\item $V_t= \lim_{s \to t} V_s$ (lower
limit!!)
\end{enumerate}

\begin{theorem}
\label{theorem-persi-modul-sum-integ}
Any persistence module is a sum of integral
persistence modules,
$$
(V,\Pi) \cong \bigoplus_{I \in B(V)}F(I).
$$
\end{theorem}

\begin{example}
\label{example-hart-and-sphere}
Heart and sphere. There is a noise in the
persistence module
of the heart due to an unnecessary critical
point.
\end{example}

$(M^{2n},\omega)$ a {\it Liouville domain} is
a compact
symplectic manifold and $X \in
\mathfrak{X}(M)$
with $X \cap\partial M$ (pitchfork, i.e.
transversal intersection!)
pointing outwards
and preserved by the symplectic form, i.e.
$\mathcal{L}_X \omega=\omega$
($\omega=d \alpha$).

When we restrict $\omega$ to the boundary,
we obtain a contact form and get some
interesting dynamics.

A Lagrangian $(L,\partial L) \subset
(M,\partial M)$
is {\it asympotically conical} if
\begin{enumerate}
\item $\partial L \subset \partial M$ is
Legendrian.
\item $L \cap [1-\varepsilon,1] \times
\partial M=
[1-\varepsilon,1] \times \partial L$.
\end{enumerate}

\begin{remark}
\label{remark-Hamiltonian}
Take a Hamiltonian $H: \hat{M}\to \mathbb{R}$
such that
$$
\begin{cases}
H(r,x)=h(r)\qquad &r=1 \\
H(r,x)=rT-B\qquad &
\end{cases}
$$
then $X_H=h'(r)R_\alpha$.

For $L_0,L,A\subset$ Lagrangians, $H$ linear
at infinite,
then $A_{H}^{L_0 \to L_1}$,
$$
\begin{aligned}
A_{H}^{L_0 \to L_1}: P_{L_0\to L_1}
&\longrightarrow \mathbb{R} \\
\gamma &\longmapsto \int_0^1 \gamma^* \alpha-
\int_0^1H(x(t))dt
\end{aligned}
$$
where $P_{L_0\to L} =\{\gamma:[0,1] \to
\hat{M}:\gamma(0) \in L_0,
\gamma(1) \in L_1\}
$ is the set of chords.
\end{remark}

\begin{remark}
\label{remark-1-chords-are-crticial-points}
$\text{crit}(A_{H}^{L_0 \to
L_1})=\{\text{1-chords of $X_H$ from
 $L_0$ to $L_1$}\}$.
\end{remark}

Putting a metric on $P_{L_0 \to L_1}$ we can
consider $\varphi:\mathbb{R} \times
[0,1] \to \hat{M}$, solutions of some PDE
which is some kind
of generalization of a gradient, $- \nabla
A_{H}^{L_0 \to L_1}$.
These solutions can be put in a moduli space
$$
\tilde{\mathcal{M}}(x_-,x_+,H,J)=
\{\varphi\text{
solutions s.t. …}\}
$$
Then we define a boundary operator
$\partial$.
\begin{theorem}
\label{theorem-homology}
$\partial^2=0$
\end{theorem}

So that we have a homology, called {\it
wrapped Floer homology}
$HW^t(H,L_0,L_1,J)$

\begin{remark}
\label{remark-embedding}
We have $H \leq  K \rightsquigarrow
HW^t(H,L_0,L_1,J) \to
HW^t(K,L_0,L_1,K)$.
\end{remark}

\begin{definition}
\label{definition-direct-limit}
For $t \geq 0$
$$
HW^t(M,L_0,L_1)=\underline{\lim }_H
HW^t(H,L_0,L_1,J)
$$
\end{definition}

(Where we have taken direct limit.)

Taking direct limit of the homology,
we make sure the homology theory is
independent of the
choice of objects (I think, complex structure
and Hamiltonian)
we used to construct it.


\begin{proposition}
\label{proposition-persistence-module}
$t \to HW^t(M,L_0,L_1)$ is a peristence
module $B(M,L_0,L_1).$
\end{proposition}

Finally we can define {\it barcode entropy}.
Fix $\varepsilon>0$,
$t \geq 0$,
$$
b_\varepsilon(M,L_0,L_1,t)=
\# \{\text{ of bars in }B(M,L_0,L_1)
\text{with length $\geq \varepsilon$
and start before $t$}\}
$$
Then
$$
\bar{h}^{HW}(M_0,L_0,L_1)=
\lim_{\varepsilon \to 0} \text{lim sup}_{t
\to \infty}
\frac{\text{log}^+(b_\varepsilon(M,L_0,L_1,t)
}{t}
$$

\medskip\noindent
Consider a contact manifold
$(\Sigma,\lambda,L_0,L_1)$
and $A$ the Lagrangians. (Example raising
question of filling.)

\begin{theorem}[M '24]
\label{theorem-entro-indep-filli}
$\bar{h}^{HW}$ is independent of the filling.
\end{theorem}

\begin{theorem}[M '24]
\label{theorem-bound}
$\bar{h}^{HW}(M,L_0,L_1) \leq
h_{\text{top}}(\alpha)$.
\end{theorem}

\begin{theorem}[M '25]
\label{theorem-restriction}
Consider $(M,L_0,L_1)$.
Let $K$ be a compact topologically transitive
hyperbolic set
for the Reeb flow $\alpha$. Assume $W_\delta
^s (q) \subset \partial L_0$,
$W_\delta^s(p) \subset \partial L_1$. Then
$$
\bar{h}^{HW}(M,L_0,L_1)\geq
h_{\text{top}}(\alpha |_{K})>0.
$$
\end{theorem}
Which says it captures dynamics beyond
unconditional phenomena.
In lower dimensions these tend to coincide,
but in higher dimension we don't
know. This is related to
 ``the sup over hyperbolic sets […]''

Here's a conjecture:
$$
\text{sup}_{L_0,L_1}\bar{h}^{HW}(M,L_0,L_1)=
h_{\text{top}}(\alpha).
$$
\medskip\noindent
{\bf Extra comments.}
One of the aims is to describe topological
entropy $h_\text{top}$
using Floer theory. Theorems by
Çineli-Ginzburg-Gürel show bounds
of topological entropy and barcode entropy
(one of which is for
hyperbolic sets).

There is a notion of  {\it admissible
Hamiltonian and Reeb vector fields}
which is related to some asymptotical
behaviour  ``linear at infinity''.
I understand that admissible vector fields
give
the interesting chords for the Floer homology
construction.

\section{Revisiting cotangent bundles}
\label{section-revisiting-cotangent-bundles}

\noindent
Mieugel Cueca, KU Leuven.
Symplectic Geometry Joint Seminar, IMPA.
September 5, 2025.

\medskip\noindent
{\bf Abstract.} Cotangent bundles provide key
examples of symplectic
manifolds. On the other hand, one can think
of Lie groupoids as generalizations
of manifolds. In this context, Alan Weinstein
constructed their cotangent
bundles and proved that they are so-called
symplectic groupoids. In this talk, I
will recall this construction and explain
what happens when one replaces a Lie
groupoid with a Lie 2 (or n)-groupoid. If
time permits, I will exhibit some of
their main applications. This is joint work
with Stefano Ronchi.

\medskip\noindent

Recall the basic properties of the cotangent
bundle $T^*M$ for a symplectic manifold:
\begin{enumerate}
\item It's a vector bundle.
\item $\left\langle{}\cdot,\cdot\right\rangle{}:TM \otimes
T^*M \to \mathbb{R}_M$
is the dual pairing.
\item $\omega_{\text{can}}\in \Omega^2(T^*M)$
is symplectic.
\item $\mathcal{L}_\varepsilon
\omega_{\text{can}}=\omega_{\text{can}}$,
$\varepsilon$ Euler vector field.
\end{enumerate}

\medskip\noindent
{\bf Main goal.} Reproduce the above for Lie
$n$-groupoids.

For $n=0$ we get the above situation. For
$n=1$ [Duzud-Weinstein], [Prezun],
for $n\geq 2$ ? and we care about $n=2$.

\begin{definition}
\label{definition-groupoid}
A {\it Lie $n$-groupoid}
$\mathcal{G}:\Delta^{\text{op}}\to
\mathsf{Man}$
such that
$$
P_{e,j}:\mathcal{G}_{\ell}\to \Lambda_j^\ell
\mathcal{G}
$$
are surjective submersions $\forall  \ell,j$
are differomorphisms for $\ell > n$.
\end{definition}

\begin{remark}
\label{remark-generators}
This sort of manifolds-valued presheaf
category is generated by
$$
\begin{aligned}
d_i^\ell:\mathcal{G}_\ell \to
\mathcal{G}_{\ell-1}&\text{face maps}
0 \leq  i,j \leq  \ell\\
s_j^\ell:\mathcal{G}_\ell \to
\mathcal{G}_{\ell+1}&  \qquad
\text{degeneracies}
\end{aligned}
$$
\end{remark}

\medskip\noindent
The tangent is a functor, it satisfies
$$
T_\bullet(\mathcal{G})=T_k(\mathcal{G})=
T\mathcal{G}_k
$$
(It looks like $T$ preserves diagrams.)

\medskip\noindent
{\bf Dold-Kan.} The category $\mathsf{SVect}$
of simplicial vector spaces
has objects
$$
\xymatrix{
\mathbb{V}_\bullet&
\mathbb{V}_n\ar[r]&\cdots
\ar[r]&\mathbb{V}_2\ar[r]^{\text{3 arrows}}&
\mathbb{V}_1\ar[r]^{\text{2
arrows}}&\mathbb{V}_0
}
$$
where the $\mathbb{V}_i$ are vector spaces.

There is a functor
$$
\mathsf{SVect}\overset{N}{\to}\{\text{chain
complexes}\geq 0\}
$$
$$
\mathbb{V}_{\bullet}\to N\mathbb{V}=
\left(N_\ell \mathbb{V}\Ker
P_{\ell,\ell},\partial=d_\ell\right)
$$

\begin{theorem}[Dold-Kan]
\label{theorem-Dold-Kan}
That's an equivalence of categories. [Confirm
this!]
\end{theorem}

Those categories are monodial:
$$
(\mathsf{SVect},\otimes),\qquad
(\mathbb{V}_\bullet \otimes \mathbb{W})_\ell
=\mathbb{V}_\ell \otimes \mathbb{W}_\ell
$$
$$
(\mathsf{ch}_{\geq 0},\otimes),\qquad (V
\otimes W)_i=
\bigoplus_{\ell +k=1}V_\ell \otimes W_k
$$
And $N$ is Lax monoidal with Lax structure
given by the
Eilenberg-Zilber map, though we won't explain
the details of this.

There are duals given by internal Hom:
$$
\begin{aligned}
\mathbb{V}^{n*}&=\underline{\Hom}(\mathbb{V},
B^n\mathbb{R})
\end{aligned}
$$
Where the internal Hom is given by
$$
\underline{\Hom}(\mathbb{V},B^n\mathbb{R})
_\ell
=\Hom_{\mathsf{SVect}}(\mathbb{V}
\otimes\Delta_n[\ell],B^n\mathbb{R})
$$
for an object
$\Delta[\ell]=\mathbb{R}[\Delta[\ell]]$.


\medskip\noindent
{\bf Properties.}
\begin{enumerate}
\item $\mathbb{V}^{n*}$ isa simplicial vector
bundle.
\item $N(V^{n*})$ and $N(\mathbb{V})^*[n]$ is
a quasi isomorphism.
\item
$\left\langle{}\cdot,\cdot\right\rangle{}:\mathbb{V}\otimes
\mathbb{V}^{n*}\to
B^n\mathbb{R}$ is non-degenerate on homology.
\item $\mathbb{V}\hookrightarrow
(\mathbb{V}^{n*})^{n*}$
Mont. a eq.
\end{enumerate}

\medskip\noindent
{\bf The vector bundle case.}
$$
\text{Maps}(\Delta[i],\mathcal{G})_k=
\Hom_{\mathsf{SSet}}(\Delta[i]\times\Delta[k]
,\mathcal{G})
$$
\begin{proposition}
\label{proposition-on-an-n-Lie-groupoid}
Let $\mathcal{G}$ be a Lie $n$-groupoid.
\begin{enumerate}
\item $\text{Maps}(\Delta[i],\mathcal{G})$
Lie $n$-groupoids ME $\mathcal{G}$.
\item
$\text{Maps}(\Delta[i],\mathcal{G})_0=
\mathcal{G}_i$.
\item $\text{ev}:\Delta[i]\times
\text{Maps}(\Delta[i],\mathcal{G})\to
\mathcal{G}$.
\end{enumerate}
\end{proposition}

[Staircase looking diagram.]

\begin{definition}
\label{definition-what-is-this}
$\mathcal{G}_\bullet$.
$$
T_i^{n*}\mathcal{G}=\Hom_{\mathsf{SVect}}(1^*
\Pi_{\Delta[i]}(T\mathcal{G}),
B^n\mathbb{R}_{\mathcal{G}_i})
$$
$$
(d_j,F)_{K|d_j\mathcal{G}}(x^a)=(F_k)
|_{\mathcal{G}}(x^{\delta_ja}).
$$
\end{definition}

\begin{proposition}
\label{proposition-properties-of-that}
$\mathcal{G}$ Lie $n$-groupoid, then
$T^{n*}\mathcal{G}$ satisfy
\begin{enumerate}
\item is a vector bundle $n$-groupoid
\item dual to $T\mathcal{G}$
$$
\left\langle{}\cdot,\cdot\right\rangle{}:T\mathcal{G} \otimes
T^{n*}\mathcal{G}
\to B^n\mathbb{R}_{\mathcal{G}}
$$
non-degenerate on homology.
\item $n$-shifted symplectic
$$
T^{n*}_n\mathcal{G}\overset{p}{\to}
T^*\mathcal{G}_n
$$
and $p^*\omega_{\text{can}}$.
\end{enumerate}
\end{proposition}

[More computations I missed]

\section{Holomorphic extensions of s-proper
Lie groupoids}
\label{section-holom-exten-s-prope-lie}

\noindent
Rui L. Fernandes, Instituto Superior Técnico
(Lisboa)/University of Illinois
at Urbana-Champaign.
Symplectic Geometry Seminar, IMPA.
September 17, 2025.

\medskip\noindent
{\bf Abstract.} Every smooth manifold admits
a compatible analytic
structure, and a classical result of
Whitney–Bruhat states that any analytic
manifold has a holomorphic extension. Lie
groups also admit compatible analytic
structures, and another classical result, due
to C. Chevalley, shows that any
compact Lie group has a holomorphic extension
to a complex Lie group. D.
Martínez Torres has shown that any proper Lie
groupoid admits a compatible
analytic structure. I will discuss an
extension of the classical results of
Whitney–Bruhat and Chevalley, establishing
that any s-proper Lie groupoid has a
holomorphic extension. This talk is based on
recent joint work with Ning Jiang
(arXiv:2508.18036).

\medskip\noindent

\begin{theorem}[Whitney-Beuhat]
\label{theorem-Whitney-Beuhat}
If $M$ is analytic there exists a complex
manifold  $M_\mathbb{C}$
together with an analytic map $i: M \to
M_\mathbb{C}$
which is totally real ($T_HM_{\mathbb{C}}=TM
\oplus J(TM)$)
and
\begin{enumerate}
\item For every complex manifold $X$ and
every
analytic map $\phi:N \to X$ there exists
$\phi^*:U \to X$
holomorphic contained open $M \subset U
\subset M_\mathbb{C}$.
\item If $\psi:V \to X$ holomorphic on $M
\subset V \subset M_\mathbb{C}$ then
$\psi=\phi^*$ where $\phi=\psi \circ i$ on a
possibly
smaller open contained in $U \cap V$.
\end{enumerate}
\end{theorem}

\begin{theorem}[Chevalley]
\label{theorem-Chevalley}
If $G$ is a Lie group, there exists a complex
Lie group $G_{\mathbb{C}}$ and a morphism
$i:G \to G_\mathbb{C}$
satisfying the following universal property.
For every complex Lie group $H$ and morphism
$\phi:G \to H$ there exists a unique
holomorphic map $\phi^*:G_\mathbb{C} \to H$
such that
$\phi=\phi^* \circ i$.
\end{theorem}

Here's the construction of $G_\mathbb{C}$:

$$
\xymatrix{
N\subset \tilde{G}\ar[r]\ar[d]&G^*\ar[d]\\
\tilde{G}/N \simeq
G\ar[r]_{i}&G_\mathbb{C}=G^* /\overline{i^*N}
}
$$
where $\tilde{G}$ is the universal cover
group of $G$,
 $G^*$ is the group that integrates to
the complexification of the Lie algebra of
$G$,
i.e.
$\text{Lie}(G^*)=\mathfrak{g}_\mathbb{C}$
and $\overline{i^*(N)}$ is the smallest
closed
normal complex Lie subgroup of $G^*$
containing $i^*(N)$.

\begin{example}
\label{example-complexification-groupoid}
$$
\xymatrix{
\widehat{\text{SL}_2(\mathbb{R})}\ar[d]
\ar[dr]\\
\text{SL}_2(\mathbb{R})\ar@{^{(}->}[r]&
\text{SL}_2(\mathbb{C})
}
$$
which is a simple case,
but consider instead
$\tilde{\text{SL}_2}(\mathbb{R})\times
\mathbb{R}$
and the normal subgroup
$N=\left\langle{}a\right\rangle{}\times \left\langle{}\lambda\right\rangle{}$
for irrational $\lambda$.
Then we obtain
$$
\xymatrix{
\widehat{\text{SL}_2(\mathbb{R})}
\times\mathbb{R}\ar[r]\ar[d]
&\text{SL}_2(\mathbb{C})\times\mathbb{C}
\ar[d]\\
\widehat{\text{SL}_2(\mathbb{R})}
\times\mathbb{R}/N\simeq
G\ar[r]&
G_\mathbb{C}
}
$$
and $G_\mathbb{C}$ is 3 complex dimensions!
It is not
$\text{SL}_2(\mathbb{C})\times\mathbb{C}$.
\end{example}

\section{Lagrangian blow-up and blow-down for
4-dimensional symplectic manifolds
}
\label{section-blow-up-down}

\noindent
Misha Verbitsky, IMPA.
Geoemtric Structures Seminar, IMPA.
October 23, 2025.

\medskip\noindent
{\bf Abstract.} The usual (complex geometric)
blow-up has a symplectic
version, called the symplectic cut. I will
introduce a variant of this
construction, which is valid only in
symplectic category, called ``Lagrangian
blow-up'', and its inverse, called Lagrangian
blow-down. Lagrangian blow-down
takes a Lagrangian sphere in a 4-dimensional
symplectic manifold and contracts
it to a symplectic orbifold with a double
point. Given a symplectic orbifold
with a double point, Lagrangian blow-up
produces a symplectic manifold, and this
construction is inverse to the Lagrangian
blow-down. I will explain why these
constructions are functorial, that is,
defined on the corresponding symplectic
Teichmuller spaces. This is used to prove an
orbifold counterexample to a famous
conjecture, sometimes attributed to
Donaldson, who asked whether any symplectic
form on a K3 surface is compatible with a
Kahler structure. I will prove that
this is false for orbifolds: some K3
orbifolds admit symplectic forms not
compatible with any Kahler structure. The
same argument is used to produce a
countable family of Lagrangian spheres in a
K3 surface which are not Lagrangian
isotopic to special Lagrangian spheres. This
is a joint work with Michael Entov.



\medskip\noindent


Consider the total space of the cotangent
bundle of $\mathbb{C}P^1$.
It's $\mathcal{O}(-2)$.
Let $\gamma \in H^{0}(\mathcal{O}(2))$.
$\gamma$ defines a fiberwise linear function
on $T^*\mathbb{C}P$.

The blow-up of $\mathbb{C}^2/\pm 1$, which is
the orbifold $\mathbb{C}^2$ with
 a doble point, is the cotangent bundle $T^*
 \mathbb{C}P^1$.
That is,
$T^*\mathbb{C}P^1 \to \mathbb{C}^2/\pm 1$.
This is called {\it crepant resolution}.
Consider the $dx \wedge dy$, which is
holomorphically symplectic
on $\mathbb{C}^2$ and descends to the
quotient $\mathbb{C}^2/\pm 1$
as $\Omega_0$.

\begin{proposition}
\label{proposition-}
$\pi:T^*\mathbb{C}P^1 \to \mathbb{C}^2/\pm
1$.
Then $\pi^*\Omega_0$ is holomorphically
symplectic.
\end{proposition}

\begin{proof}
Actually, it gives the standard symplectic
form on $T^*\mathbb{C}^2$.
\end{proof}

\medskip\noindent
Now we shall define the Langrangian blow-up.
It's a procedure that starts with a
symplectic manifold
with double points and outputs a symplectic
manifold
without double points.

\begin{definition}[Lagrangian blow up]
\label{definition-lagrangian-blow-up}
Let $(B,\omega)$ be the standard ball in
$\mathbb{R}^n$,
$dz_1 \wedge dz_2 +dz_3 \wedge dz_4$.
Consider the symplectic orbifold
$(B,\omega)/\pm 1$.
Consider $M$ be a symplectic orbifold
with double point singularities.
Darboux coordinates in a neighbourhood of $0$
with $(B,\omega)/\pm 1$ resolve
singularities.
\end{definition}

This definition has the caveat that there ar
e lots of choices involved.
We shall later that these choices do not
alter the result,
namely Theorem
\ref{theorem-lagrangian-blow-up-well-defined}
.

To define the Lagrangian blow-down we first
recall

\begin{theorem}[Weinstein's Lagrangian
neighbourhood]
\label{theorem-weins-lagra-neigh}
Let $X \subset M$ be a compact Lagrangian
submanifold in $(M,\omega)$.
Then there exists a neighbourhood $U$ of $X
\subset M$
which is symplectomorphic to a neighbourhood
of $X$
in $X \subset T^*M$.
\end{theorem}

\noindent
Let $S \subset M$ be a Lagrangian sphere
(diffeomorphic to sphere
and Lagrangian).

\begin{definition}
\label{definition-lagrangian-blow-down}
The {\it Lagrangian blow-down} is obtained
from $M$
by removing $W \subset M$ and gluing $W_0$ in
its place,
where $W_0$ is such that

$$
\xymatrix{
W
\ar@{^{(}->}[r]\ar[d]&T^*\mathbb{C}P^1\ar[d]
\\
W_0\ar@{^{(}->}[r]&\mathbb{C}^2/\pm 1.
}
$$
\end{definition}

\medskip\noindent
Now we recall Teichmüller spaces.

Let $\Gamma(\Lambda^2M)$ be the space of all
2-forms on a manifold $M$.
The space of symplectic forms is a subspace
of the closed 2-forms on $M$,
which in turn is a subspace of
$\Gamma(\Lambda^2M)$.
Equip this space with the $C^\infty$
topology.
The {\it Teichmüller space} is
$\text{Symp}/\text{Diff}_0$.

Two symplectic structures are called {\it
isotopic} if they lie
in the same orbit of $\text{Diff}_0$.
The {\it period map}
$\text{Per}:\text{Teich}_s \to
H^{2}(M,\mathbb{R})$
maps a symplectic structure to its cohomology
class.

\begin{theorem}[Moser, 1965]
\label{theorem-period-map-is-local-diffeo}
The period map is a local diffeomorphism.
\end{theorem}

\noindent
This implies that $\text{Teich}_s$ is smooth.

\begin{theorem}[Moser's trick]
\label{theorem-mosers-trick}
Let $\omega_t$, $t \in S$ be a smooth family
of symplectic
structures, parametrized by a connected
manifold $S$.
Assume that the cohomology class
$[\omega_t]$ is constant.
Then all $\omega_t$ are diffeomorphic.
\end{theorem}

\noindent
This means that the fibers of $\text{Symp}\to
H^{2}(M)$
are the orbits (of the action of
$\text{Diff}_0$, I suppose).

\medskip\noindent
Consider a Lagrangian submanifold $L \subset
(M,\omega)$.
Let us assume for simplicity (though the
results
can also be obtained without this) that
$b_1(L)=0$,
i.e. first cohomology group is zero.

Let $\text{Symp}_L$ be the set of symplectic
forms vanishing on $L$,
and $\text{Diff}_0^L$ the connected component
of diffeomorphisms
preserving $L$.
Define
$\text{Teich}_L=\text{Symp}_L/\text{Diff}
_0^L$.

\begin{theorem}
\label{theorem-}
Let $V \subset H^{2}(M,\mathbb{R})$ be the
space of all
cohomology classes on $M$ which vanish on
$L$,
and $\text{Per}_L:
\text{Symp}_L/\text{Diff}^L_0 \to V$
take $(M,\omega)$ to the cohomology class of
$\omega$.
Then Per is a local diffeomorphism.
\end{theorem}

\begin{proof}
We used a version of Moser's trick: for
$\omega_t$ a family
of symplectic forms with $\omega_t|_{L}=0$,
with $[\omega_t]$ constant, then $\omega_t$
are $\text{Diff}_0^L$-isotropic.
\end{proof}

\noindent
As a corollary we obtain that this
Teichmüller space
classifies Lagrangian submanifolds:

\begin{lemma}
\label{lemma-techm-class-lagra-subma}
Assuma that $L$ and $L'$ are Lagrangian
submanifolds in $(M,\omega)$,
and $\varphi \in \text{Diff}_0(M)$ a smooth
isotopy such that
$\varphi(L)=L'$. Then $(\omega,L)$ and
$(\varphi^*\omega,L)$
represent the same point in $\text{Teich}_L$
if and only if
$L$ and $L'$ are Lagrangian isotopic in
$(M,\omega)$.
\end{lemma}

\noindent
Again: points of this Teichmüller space are
isotopy classes
of Lagrangian submanifolds.

\medskip\noindent
\begin{theorem}
\label{theorem-blow-down-diffeomorphism}
Let $L$ be a Lagrangian 2-sphere in a K3
surface  $M$
and $\hat{M}$ the Lagrangian blow-down of
$M$.
The blow-up/blow-down
is a diffeomorphism $\text{Teich}_L(M) \to
\text{Teich}(\hat{M})$.
\end{theorem}

\medskip\noindent
Now we prove the theorem which asserts that
the Lagrangian
blow-up is independent of choices.

\begin{theorem}
\label{theorem-lagrangian-blow-up-well-defined}
Let $(\hat{M},\hat{\omega})$ be a orbifold.
Then its blow-up
$(M,L)$ defines a point in $\text{Teich}_L$
independent of the choices.
\end{theorem}

\begin{proof}
By taking local Darboux coordinates, and then
Moser's trick.
\end{proof}

\noindent
Actually, the same happens for blow-downs

\begin{theorem}
\label{theorem-same-for-blow-down}
Let $(M,L,\omega)$ be a 4-manifold with a
Lagrangian sphere $L$, and
$(\hat{M},\hat{\omega)}$ be its Lagrangian
blow-down.

Then the corresponding point in
$\text{Teich}_{\hat{M}}$ is
independent from the choices made.
\end{theorem}

\begin{proof}
Similar.
\end{proof}

\medskip\noindent
A conjecture by Donaldson ways that all
symplectic structures on a K3
surface are compatible with a Kähler
structure. This conjecture
(although it's commonly believed to be true)
is still open.
When trying to prove that this would hold for
orbifolds,
Misha and collaborators realised it's
actually false.

Note: Seiberg-Witten invariants may be used
to tell whether a symplectic
form is compatible with a complex structure.

\begin{theorem}
\label{theorem-conje-false-orbif}
There exists an orbifold symplectic K3
surface $M$ with a single
double point not admitting an orbifold Kähler
structures.
\end{theorem}

\begin{proof}
\noindent
Step 1. Consider
$$
\text{Teich}_{\text{Kähler-symplectic}}
(\hat{M})
=\{\eta \in H^{2}(\hat{M}):\eta^2>0\}.
$$
As follows from [Amerik-V., 2005], the
Teichmüller space of symplectic structures
compatible
with a hyperkähler structure is Hausdorff and
connected
(the same argument works for Kähler K3
orbifolds).

Essentially, if there wasn't any orbifold
as required, the Teichmüller space
classifying
Lagrangian submanifolds would be connected,
which
is impossible by [Seidel 2000].
\end{proof}

\noindent
In fact, Seidel's result is essentially
mirror symmetry.

\medskip\noindent
Now we discuss special Lagrangian
submanifolds.

\begin{definition}
\label{definition-phase}
Let $(M,I,\omega,\Omega)$ be a Calabi-Yau
manifold, where
$\Omega \in \Lambda^{n,0}_I(M)$ is
nondegenerate and $\omega$ is Kähler.
Let $L$ be a Lagrangian submanifold.
Define the {\it phase} as
$$
\begin{aligned}
\text{phase}: L &\longrightarrow \Lambda^nL
\otimes \mathbb{C} \\
x &\longmapsto \Omega|_{T_xL}.
\end{aligned}
$$
\end{definition}

\noindent
A better definition is:
$$
\begin{aligned}
\text{phase}: L &\longrightarrow
S^1=\text{U}(1) \\
x &\longmapsto \frac{\Omega|_{T_xL}}{|\Omega
T_xL|}.
\end{aligned}
$$

\begin{definition}
\label{definition-special-lagrangian}
A special Lagrangian submanifold is the phase
is constant.
\end{definition}

\noindent
Special Lagrangian submanifolds have
deformation space of
dimension 1, it's unobstructed. They are
calibrated, so minimal.

\begin{theorem}
\label{theorem-lagrangians-of-k3}
$L \subset$ K3 is Lagrangian is isotopic to a
special Lagrangian
if and only if its blow-down is Kähler type.
\end{theorem}

\medskip\noindent
One of the basic results about special
Lagrangian
submanifolds is

\begin{lemma}[Hitchin]
\label{lemma-hitchin}
If $(M,\Omega)$ is holomorphically symplectic
and $X \subset M$
is Lagrangian with respect to
$\text{Re}\Omega$ and $\text{Im}\Omega$
then $X$ is complex.
\end{lemma}

\noindent
As corollaries:

\begin{lemma}
\label{lemma-corollary-hitchin}
Let $(M,I,J,K,g)$ be a hyperkähler manifold
of real dimension $4n$,
considered as a Kähler manifold
$(M,I,\omega_I)$ and $\phi:\Omega^n \in
\Lambda^{2n,0}(M,I)$ the corresponding
holomorphic volume form.
Consider a Lagrangian submaniold $S \subset
(M,\omega_I)$
such that $a \omega_J + b\omega_K|_S=0$ for
some nonzero real numbers
$a,b \in \mathbb{R}$ such that $a^2+b^2=1$.
Then $S$ is special Lagrangian, and,
moreover,
it is holomorphic with respect to the complex
structure
$-a K+b J$.
\end{lemma}

\begin{lemma}
\label{lemma-corollary2-hitchin}
Let $\omega_I,\omega_J$ and $\omega_K$ on
$M$, with $\dim M=4$.
Let $S \subset (M,\omega_I)$ be Lagrangian.
Then $S$ is special Lagrangian if and only if
$S$ is complex analytic
on $a J + b K$ for two nonzero numbers  $a,b
\in \mathbb{R}$ such that
$a^2+b^2=1$.
\end{lemma}

\noindent
Now an auxiliary claim:

\begin{lemma}
\label{lemma-k3-type}
$\Omega$ is holomorphic symplectic
\end{lemma}

\begin{theorem}
\label{theorem-}
Let $S\subset M$ be a Lagrangian sphere in a
K3 surface.
Then  $S$ is isotropic to a special
Lagrangian if and only if
$M$ is of Kähler type.
\end{theorem}

\section{Systolic inequality for tori}
\label{section-systolic-inequality}

\noindent
Dimitri Korshunov, IMPA.
Minimal Surfaces end of course presentation.
October 24, 2025.

\medskip\noindent
This is due to L. Guth.

\begin{definition}
\label{definition-systole}
Given a Riemannian manifold,
a {\it systole} $\text{Sys}(M,g)$ is the
infimum of the lengths
of noncontractible paths.
\end{definition}

\begin{theorem}[Gromov]
\label{theorem-gromov}
For any metric on a torus $T^n$,
$$
\text{Sys}(T^n,g) \leq C_n
\text{Vol}(T^n,g)^{\frac{1}{n}},
$$
where $C_n$ is a constant depending only on
the dimension
of the torus $n$.
\end{theorem}

\noindent
This also holds for $\mathbb{R}P^n$. We will
use induction,
for which we need the following stronger
statement:
given a number bounded by half the systole,
there exists a ball
which is {\it big}:

\begin{theorem}[Guth]
\label{theorem-guth}
For $(T^n,g)$, given
$R<\frac{1}{2}\text{Sys}(T^n,g)$,
then there exists a point $p$ on $T^n$
such that the volume of ball of radius $r$
with center in $p$
satisfies
$$
\text{Vol}B(p,R) \geq C_nR^n.
$$
\end{theorem}

\noindent
It was conjectured by Gromov that $C_n=1$.

\medskip\noindent
Before stating the theorem that we will
prove, we introduce the
following definition that will allow us to
measure cohomological cycles:

\begin{definition}
\label{definition-measu-cohom-cycle}
Let $\alpha \in H^{1}(M_g,\mathbb{Z}_2)$.
$$
L(\alpha)
=\text{inf}\{|\lambda|:\lambda\text{
1-cycles, }\alpha(\lambda)\neq 0\}.
$$
\end{definition}

\noindent
The theorem that we will prove, which implies
Theorem \ref{theorem-guth}
is the following

\begin{theorem}
\label{theorem-bound-classes}
There exists a constant $\varepsilon_n>0$
such that given a manifold $(M,g)$, and
$\alpha_1,\ldots,\alpha_n \in
H^1(M,\mathbb{Z}_2)$,
\begin{enumerate}
\item $L(\alpha_i)>2R$
\label{item-positive}
\item
\label{item-cup} $\alpha_1 \cup  \alpha_2
\cup \ldots \cup \alpha_n \neq 0 \in
H^n(M,\mathbb{Z})$.
\end{enumerate}

\noindent
Then there exists $p \in M$ such that
$$
\text{Vol}B(p,R) \geq \varepsilon_nR^n.
$$
\end{theorem}

\noindent
In particular $T^n$ and $\mathbb{R}P^n$
satisfy (\ref{item-cup}).

Before we proceed to the proof we need the
following definition.
Essentially we ask that the surface minimizes
volume in its homology
class up to some constant $\delta$.

\begin{definition}
\label{definition-delta-minimal-surface}
A  {\it $\delta$-minimal surface} is an
embedded hypersurface $Z$
such that for any $Z'$ such that $[Z'-Z] = 0$
in  $H_{n-1}(M,\mathbb{Z}_2)$,
then  $\text{Vol}Z < \text{Vol}Z'+\delta$.
\end{definition}

\noindent
Moreover, we need the following

\begin{lemma}[Stability]
\label{lemma-stability}
Let $(M,g)$ de closed Riemannian manifold and
$\alpha \in H^1(M,\mathbb{Z}_2)$.
Suppose $Z$ is Poincaré dual to $\alpha$ and
that $Z$ is $\delta$-minimal.
If $R<\frac{1}{2}L(\alpha$) then the
following
holds for any $p \in M$:
$$
|Z \cap \text{Vol}B(p,\frac{R}{2}|\leq
\frac{2}{R}|\text{Vol}B(p,R)|+\delta.
$$
\end{lemma}

\begin{proof}[Proof of the stability lemma]
\medskip\noindent
Step 1. Using commas formula, namely
$$
\text{Vol}(B_R)=\int_0^R S_tdt,
$$
we see that there exists $t$ such that
$\frac{R}{2}<t<R$ and
$$
|S(p,t)|_R< \frac{2}{R}|\text{Vol}B(p,R)|.
$$

\medskip\noindent
Step 2. Consider
$[Z \cap B(p,t)] \in
H^{n-1}(B(p,t),S(p,t),\mathbb{Z})$.
This can be interpreted geometrically since
we are using relative homology.

\medskip\noindent
Step 3. We claim that $[Z \cap B(p,t)]=0$.
By contradiction,
if it weren't trivial, then there would exist
a cycle $H_1(B,\mathbb{Z})$
such that $Z' \cap \gamma \neq 0$ since
Poincaré pairing is nondegenerate ---
namely any cocycle that is not zero pairs
nontrivially with some cycle.
Then we will just split the cocycle in very
small parts,
and we will reach a contradiction with the
fact that
$R<\frac{1}{2}L(\alpha)$.

We will state this as another lemma.
Essentially it says that any
cycle in homology can be decomposed as a sum
of cycles
with small lengths.

\begin{lemma}[Curve factoring, Gromov]
\label{lemma-curve-factoring-gromov}
Let $(M,g)$ be a Riemannian manifold and
$\gamma$ a cycle in $B(p,t)$.
Then $[\gamma]=\sum \sigma_i$ is such that
 $L(\sigma_i)<2t+a$ for all $i$.
\end{lemma}

\begin{proof}
We spit a cycle (a curve) into some quantity
of intervals
such that each has length smaller than $a$,
then add paths joining the center point of
length $t$,
so we count $t$ twice: to reach the point and
to get away from it.
\end{proof}

\noindent
We apply this to $\gamma$ and reach a
contradiction
with the hypothesis that
$R<\frac{1}{2}L(\alpha)$ making $a$ smaller
and
smaller.

\medskip\noindent
Step 4. Finally we use the minimality of $Z$
and the property
derived from Commas lemma to arrive at a
contradiction.
\end{proof}

\noindent
Now we prove Theorem
\ref{theorem-bound-classes}:

\begin{proof}
By induction. The case $n=1$ is trivial.
Now take the last one of our classes,
$\alpha_n$.
Let $Z$ be Poincaré dual of $\alpha_n$ and
$\delta$-minimal.
Notice that $\alpha_1\cup \ldots \cup
\alpha_{n-1}\neq 0$
since if it were, then $\alpha_1\cup
\ldots\cup \alpha_n$ would also vanish.

By induction hypothesis we get a point $p \in
Z$
such that $\text{Vol} B_Z(p,\frac{R}{2}) >
R^{n-1}$.
This means that , $|Z \cap
B(p,\frac{R}{2})|>\varepsilon_n(R/2)^{n-1}$.

By the Stability Lemma \ref{lemma-stability},
$$
\begin{aligned}
\frac{2}{R}|B(p,R)|&>\varepsilon_n\frac{R^{n-
1}}{2^{n-1}}-\delta\\
\implies
|B(p,B)|&>\frac{\varepsilon_n}{2^n}R^n-
\frac{\delta
R}{2}
> \frac{\varepsilon_n}{2^n}R^n.
\end{aligned}
$$
\end{proof}

\section{From higher Lie groupoids to higher
Lie algebroids}
\label{section-highe-group-highe-algeb}

\noindent
Matias del Hoyo, UFRJ.
Symplectic Geometry Seminar, IMPA.
October 29, 2025.

\medskip\noindent
{\bf Abstract.}
This is a report of an ongoing collaboration
with A. Cabrera, where we develop a
higher Lie theory. Starting with simple,
concrete examples, we present the
differentiation functor in a geometric,
explicit way, using simplicial methods
and differential graded algebras. We describe
how our construction extends the
classical Lie functors, relate our project
with previous works, and show a van
Est theorem linking the cohomology at the
global and infinitesimal levels.

\medskip\noindent

Upshot: Integration-differentiation relation
between Lie groups and Lie algebras
in higher versions. Groupoids are the
counterparts of Poisson manifolds, and in
the next level there are 2-groupoids and
Courant algebroids. How to integrate a
2-groupoid to a Courant algebroid? This is an
open question which motivated this
work. See arXiv
\href{http://arxiv.org/abs/2309.14105}%
{2309.14105}.

\medskip\noindent
Let $M$ be a smooth manifold.

\begin{remark}
\label{remark-smoot-funct-ring-fully-faith}
$M \mapsto C^\infty M$ is fully faithful.
\end{remark}

\noindent
Which allows us to think of manifolds as
algebraic objects.
Notice that
$\Omega^kM=\{k\text{-forms}\}
=\Lambda^k_{C^\infty M}\text{Der}(C^\infty
M,C^\infty M)$,
which is (algebraically) locally spanned
by the exact 1-forms.

Recall that the exterior differential is
the only operator $d:\Omega^kM \to
\Omega^{k+1}M$
such that $d(\alpha \wedge \beta)=d\alpha
\wedge \beta +
(-1)^{|\alpha|}\alpha \wedge d\beta$
and that $df(X)=Xf$.

This makes $(\Omega M,\wedge,d)$ a
differential graded algebra.
Recall that $\text{dom}(X,Y)=X\omega(Y)-Y
\omega(X)-\omega([X,Y])$.
We also have de Rham cohomology.

Consider a submanifold $S \subset M$. Recall
that $\mathcal{I}_S/\mathcal{I}^2_S$ is
essentially
the sections of the normal bundle.
In particular, for $M \subset M \times M$
we can see that the embedding vanishes on the
diagonal,
so there is no linear part in the Taylor
expansion,
and the second term is in fact the de Rham
differential
$df=\pi_1^*f-\pi_2^*f \text{ mod
}\mathcal{I}^2_{\Delta(M)}$.

In fact, this construction can be extended to
a simplicial manifold
(which I think is just a presheaf valued on
category of manifolds)
using projections and diagonals,
$$
\xymatrix{
\cdots & M\times M \times M
\ar @< 3pt> [r]
\ar @< 0pt> [r]
\ar @<-3pt> [r] & M \times M \ar @< 2pt> [r]
\ar @<-2pt> [r]& M.
}
$$
To this we can apply the $C^\infty$ functor
and obtain
a cosimplicial ring (not a dga?).
The idea here
is that de Rham differential is an
alternating sum of pullbacks;
as other differentials will also be.

\medskip\noindent
Now observe that for a finite-dimensional
vector space $V$,
$$
\left\{\substack{\text{Lie bracket} \\
V \times V
\xrightarrow{[\cdot,\cdot]}V}\right\}
\leftrightarrow
\left\{\substack{\text{differential on} \\
\Lambda V^*}\right\},
$$
making $\Lambda V^*$ into a differential
graded algebra.

Recall that given a Lie algebra
$\mathfrak{g}$,
we have $(\Lambda^\bullet \mathfrak{g}^*
,d)=\text{CE}(\mathfrak{g})$,
the Chevalley Eilenberg algebra,
which gives Lie algebra cohomology.

\medskip\noindent
[Missing part on Simplicial approach,
exponential map,
Simplicial manifold nerve]

\medskip\noindent
Lie algebroid: $\Delta \to M$ vector bundle
with $[\cdot,\cdot]$ on $\Gamma A$,
$A \xrightarrow{\rho}TM$ Leibniz.

\begin{remark}
\label{remark-lie-algebroids-differentials}
$$
\left\{\substack{\text{Lie algebroids on} \\
E \to M}\right\}
\leftrightarrow
\left\{\substack{\text{differentials} \\
\text{on }\Lambda\Gamma E^*}\right\}.
$$
\end{remark}

\noindent
This gives a differential graded algebra
$$
\xymatrix{
\Delta\ar[r]&M\ar[r]&\text{CE}(A)
}
$$
and
$$
\xymatrix{
C^\infty M\ar[r]^{\tilde{\rho}}&\Gamma
E^*\ar[r]^{\widetilde{[\cdot,\cdot]}}
&\Gamma E^* \wedge \Gamma E^*\ar[r]&\cdots.
}
$$

From Lie groupoids to Lie algebroids:
Lie functor/differentiation.

$$
\Omega(G)=\Gamma \Lambda^\bullet
\underbrace{\supset}
_{\text{sub dga}}
\Gamma \Lambda^\bullet T^s
\xrightarrow{-/G}\text{CE}(A_G).
$$
There's the following ecosystem around Lie
groupoids:
$$
\xymatrix{
&\text{Lie groupoids}
\ar[dl]\ar@{->>}[d]\ar[dr]^{\text{nerve}}\\
\text{Lie algebroids}&
\substack{\text{Differentiable stacks} \\
\supset \text{orbifolds}}
&\text{simplicial manifolds.}
}
$$

\medskip\noindent
Lie groupoids are models for stacks,
which are geometric objects solving
classification
problems with objects identified up to
isomorphism.
Relaxing the identification to a weaker
notion
of equivalence we must introduce higher
versions of stacks.

Moving higher,
$$
\xymatrix{
\text{Lie groupoid}\ar[r]\ar[d]_{\text{Lie}}
&\text{Simplicial manifolds}\ar[d]^?\\
\text{Lie
algebroids}\ar[r]&\text{dg-manifolds}
}
$$
dg algebra is a  dga $(A=\bigoplus_{k=1}^n
A_k \to M,\wedge,d)$
starting at $C^\infty M$, and continuing
until you get to
the sections, namely
$$
\xymatrix{
C^\infty M \ar[r]_d & A_2 \ar[r]& \cdots &
\cong \Gamma A.
}
$$
[dH-Cabrera] provide the solution to this
problem,
namely what to put in the right vertical
arrow:

\begin{theorem}[dH-Cabrera]
\label{theorem-dh-cabrera}
$$
\frac{(C^\infty_N(G),\cup,\sum(-1)A^*_i)}{J+
\delta
J}
=\text{CE}(A_G)
$$
where $A_G$ is a higher Lie algebroid
and $J$ is defined by: $f$ is in $J$ if it
has higher vanishing.
\end{theorem}


Groupoids are the counterparts of Poisson
manifolds,
and in the next level there are 2-groupoids
and Courant algebroids.
How to integrate a 2-groupoid to a Courant
algebroids?
This is an open question which motivated this
work.

\section{Hyperbolic surfaces and Mirzakhani's
curve counting}
\label{section-hyper-surfa-mizak-curve}

\noindent
Viveka Erlandson, University of Bristol, UK.
8th Brazilian School of Dynamical Systems,
IMPA.
November 4, 2025.

\medskip\noindent
{\bf Abstract.}
It’s a classical problem to study the growth
of the number of periodic orbits of
bounded length $L$ under a geodesic flow on a
manifold. When it comes to
(closed) hyperbolic surfaces Huber proved
that the number is asymptotic to $e^L
L$ and this results has since been
generalized in many directions, maybe most
famously by Margulis to negatively curved
manifolds as well as more general
flows. However, a fundamentally different
result is due to Mirzakhani who
counted the subset of those closed geodesics
on hyperbolic surfaces which have
no self-intersections (or more generally,
those inside a fixed mapping class
group orbit) and showed that in this setting
the growth is asymptotic to a
polynomial in L of degree only depending on
the topology of the surface. In this
mini-course we will give the necessary
background on hyperbolic surfaces to
understand these results, give an idea of the
proof of Mirzakhani’s theorem and
introduce some tools used such as measured
laminations and train tracks. Time
allowing we present some generalizations (for
example, to non-simple closed
geodesics or to various metrics) and
applications to the distribution of closed
geodesics.

\medskip\noindent
Our goal is to count certain closed geodesics
on hyperbolic surfaces.
By ``certain'' we main, for example, simple
or fixed topological type;
we will define these properly. Also, we may
treat more general
types of curves than only hyperbolic, e.g.
combinatorial, algebraic.
As ``closed geodesics'', we mean free
homotopy classes of curves.

As historical note,
results on coarse polynomial growth for
simple curves
were are due to Birman-Serres, Rees in the
80's.
In 2001, McShane-Runn studied asymptotic
growth
for simple curves on punctured torus.
Asymptoric growth on any hyperbolic
surface was proved by Mirzakhani for simple
curves in 2008
and then in 2016 for any topological type of
curve.

Here's an outline of this minicourse:
\begin{itemize}
\item Introduction, motivation.
\item Measured laminations, train tracks,
Thurston measure.
\item  (Geodesic) currents.
\item Curve counting theorem.
\end{itemize}

\medskip\noindent
Throughout $S$ will be a topological
(connected, orientable
of finite type, i.e. finite genus). Let
$g<\infty$ be the genus
and $r< \infty$ the number of punctures.
We assume that $\chi(S)=2-2g-r<0$,
which is equivalent to saying that $S$
admits hyperbolic metrics.
We will not consider the pair of pants,
since there's only three simple closed
geodesics on it,
making it uninteresting for our means.

Let $X$ be a hyperbolic surface homeomorphic
to $S$.
Consider the universal cover
$\tilde{X}=\mathbb{H}^2$ of $X$,
and
$\pi_1(X)=\Gamma\underbrace{<}
_{\substack{\text{discrete}
\\ \text{free}}}
\text{PSL}(2,\mathbb{R})=\text{Isom}^+
(\mathbb{H}^2)$.
Recall that a $4g$-gon yields a genus $g$
surface $S_g$ by
adequate identification of its sides.

For us, a {\it curve} $\gamma$ is a free
homotopy class
of a closed immersion $\mathbb{S}^1
\hookrightarrow S$,
which are closed $X$-geodesics.
In turn, a {\it multi-curve}
is a formal sum $\sum_{i=1}^na_i\gamma_i$
for $\gamma_i$ distinct curves and $a_i>0$,
$a_i \in \mathbb{Q}$.
The {\it length} of a curve $\gamma$ is
$\ell_X(\gamma)$ is the length of
$X$-geodesic representing $\gamma$,
and the length of a multi-curve is the sum of
the lengths
of the curves in the formal sum.

We won't use this: we can think of $X$ as
an element of Teichmüller space
$\mathcal{T}(S)\cong \mathbb{R}^{6g-6+2r}$
parametrizing complex structures on $S$.

\medskip\noindent
Notice that given $X \in\mathcal{T}(S)$,
the number of curves with length bounded by a
given number $L$,
$\#\{\gamma|\ell_X(\gamma)\leq L\}$,
is finite. This can be checked going to the
universal cover.
However, as $L \to \infty$, the number also
tends to infinity.

Huber, Hekhal (59' or later) showed that if
$X$ is hyperbolic
on $S_{g,r}$, then
$\#\{\gamma|\ell_X(\gamma)\leq L\}\sim
\frac{e^L}{L}$
asymptotically, that is,
$$
\lim_{L\to \infty}\frac{\#
\{\gamma|\ell_X(\gamma)\leq L\}}{e^L/L}=1.
$$

\medskip\noindent
What is we restrict to a fixed topological
type?

Recall that the {\it mapping class group} of
$S$ is
$\text{Map}(S):=\text{Homeo}^+(S)/
\text{homotopy}$.
Then $\gamma$ and $\gamma_0$ are of the same
(topological)
type if $\gamma \in \text{Map}(S)\gamma_0$.

We define the {\it (geometric) intersection}
of $\alpha,\beta$
distinct curves by
$$
i(\alpha,\beta)=\text{min}\{\#(\alpha' \cap
\beta'|
\substack{\alpha'\text{ homotopic to }\alpha
\\
\beta'\text{ homotopic to }\beta}\}.
$$
(In fact, this minimum is always realized by
geodesics,
but it's more a topological definition.)
The {\it self-intersection} of a curve is
$$
i(\alpha,\alpha)=\frac{1}{2}\text{min}
\{\#(\alpha' \cap
\alpha''|\alpha',\alpha''\text{
homotopic to }\alpha\}.
$$

We say $\gamma$ is {\it simple} if
$i(\gamma,\gamma)=0$.

It is known that $i(\cdot,\cdot)$ is
$\text{Map}(S)$-invariant,
and that $\{\gamma:i(\gamma,\gamma)=k\}$ is a
finite union
of topological types.

\begin{theorem}[Mirzakhani, '08, '16]
\label{theorem-mirzakhani}
Let $X$ be a hyperbolic surface, $X \cong
S_{g,r}$
and $\gamma_0$ any (multi-)curve.
Then
$$
\#\{\gamma\text{ type
}\gamma_0|\ell_X(\gamma)\leq L\}
\sim CL^{6g-6+2r},
$$
where $C>0$,
$C=\frac{c(\gamma_0)}{b_\gamma}B(X)$.
\end{theorem}

\noindent
As a corollary, we obtain
$$
\begin{aligned}
\#\{\gamma\text{ simple}|\ell_X(\gamma)\leq
L\}&
\sim\text{const.}L^{6g-6+2r}\\
\#\{\gamma|i(\gamma,\gamma)\leq
k,\ell_X(\gamma)\leq L\}
&\sim (\text{const. depending on
$k$})L^{6g-6+2r}.
\end{aligned}
$$

\medskip\noindent
As an example, consider the torus with simple
(multi-)curves.
If we put a puncture it becomes a hyperbolic
curve.
Forgetting the puncture, we have an Euclidean
flat torus,
of which we know the geodesics: they are
straight lines
with rational slope in the plane (before
quotienting).
We see that the amount of geodesics of length
bounded by $L$
is $\frac{1}{2}\mathbb{Z}^2 \cap B((0,0),L)$,
which is in
fact $\sim \pi L^2$.
A result by Mc-Shane-Rivn, 2001 shows that
 $$
\{\text{(primitive) simple curves}\}\sim
\frac{\pi}{6}L^2.
$$
Where primitive (most likely) excludes
counting
a curve twice as a multicurve.


\medskip\noindent
Let us introduce a measure on $\mathbb{R}^2$.
Let $L>0$. Then
$$
m^L=\frac{1}{L^2}\sum_{p \in
\mathbb{Z}^2}\delta_{\frac{1}{L}p}
\xrightarrow{L \to
\infty}\text{Leb}_{\mathbb{R}^2}.
$$
For $p=(a,b)$,
$\frac{1}{L}p=(\frac{1}{L}a,\frac{1}{L}b)$.
$$
\begin{aligned}
U_{\mathbb{E}}&=\{x \in \mathbb{R}^2:|x|\leq
1\}\\
m^L(U_{\mathbb{E})}&=\frac{1}{L^2}\#\{p \in
\mathbb{Z}^2
\left|\frac{1}{L}p\right|\leq 1\}\\
&=\frac{1}{L^2}\#\{p \in \mathbb{Z}^2:|p|\leq
L\}\\
&=\frac{1}{L^2}\#\mathbb{Z}^2 \cap
B_{\mathbb{R}^2}((0,0),L)\\
&\to\text{Leb}_{\mathbb{R}^2}(U_{\mathbb{E})}
=\pi.
\end{aligned}
$$

\medskip\noindent
Let's go back to study the geodesics on a
torus.
First we recall a fact, which is that
on any hyperbolic surface $S$ there exists a
compact $K \subset S$ such that all
simple closed geodesics stay in $K$.

Now consider a geodesic given by a curve of
irrational slope,
which is dense in the torus. However,
it is not dense on a hyperbolic torus
$S_{1,1}$.
We may approximate this geodesic by $q_n \in
\mathbb{Q}$ with
$q_n \to \alpha$. Each $q_n \rightsquigarrow
\gamma_n$
is a closed geodesic on $S_{1,1}$.
The lamination in this case is the Hausdorff
limit
of the $\gamma_n$.

\begin{definition}
\label{definition-lamination}
A {\it lamination} $L$ on a (hyperbolic)
surface $S$
is a compact subset of $S$ consisting of
simple, pairwise
disjoint complete geodesics.

A {\it measured lamination} $\lambda$ is
$\lambda=(L,\lambda)$ is a lamination $L$
together with a {\it transverse measure}
$\lambda$:
a measure is assigned to each arc $I$
transverse to $L$.

We write $ML(S)$ for the space of all
measured laminations on $S$.
\end{definition}

\noindent
While the definition of lamination depends on
a choice
of hyperbolic metric on the surface, the
different $ML(S)$
that we obtain with different metrics are
actually homeomorphic,
and we thus do not distinguish between them.

As an example, every simple multi-curve is a
measured lamination.

Morally we think of measured laminations as
limits of simple multicurves:
simple multi-curves are dense in $ML(S)$.

\medskip\noindent
Now we will discuss train tracks, which will
give us a combinatorial model for measured
laminations.

\begin{definition}
\label{definition-train-track}
A {\it train track} $\tau$ on $S$ is an
embedded (trivalent)
graph such that at every vertex the three
half edges
meet with the same tangent line and such that
the complementary region (cc. $S \setminus
\tau$)
are not disk
(with a certain caveat: it can be a
contractible
region as long as the vertices cannot
preserve
it's smooth incidence structure
along the homotopy),
a cyllinder (ring) nor a punctured disk.

A {\it set of weights} on $\tau$ is an
assignment
of a non-negative real number $w_e$
to each edge $e$. They satisfy the {\it
switch equations}
if at every vertex we have $w_a=w_b+w_c$
(where $b$ and $c$ ``stem'' from $a$).

We denote $w(\tau)$ the solutions to the
switch equations
and $w_{\mathbb{Z}}(\tau)$ the integral
solutions.

We say a measured lamination $\lambda \in
ML(S)$
is {\it carried} by $\tau$ ($\lambda < \tau$)
if its support can be homotopied into $\tau$.
We denote
$$
\begin{aligned}
ML(\tau)&=\{\lambda \in ML(S):\lambda <
\tau\}\\
ML_{\mathbb{Z}}(\tau)&=\{\lambda \in
ML_{\mathbb{Z}}(S):\lambda <\tau\}\\
ML_{\mathbb{Z}}(S)&=\{\lambda \in
ML(S):\substack{\text{
simple multi-curve} \\ \text{with integral
coefficients}}\}.
\end{aligned}
$$
\end{definition}

\medskip\noindent
Now we will explain the fact that there are
one-to-one correspondences
$$
\begin{aligned}
ML(\tau)&\xymatrix{\ar@{<->}[r]&}w(\tau)\\
ML_{\mathbb{Z}}(\tau)&\xymatrix{\ar@{<->}[r]&
}w_{\mathbb{Z}}(\tau).
\end{aligned}
$$

\noindent
In a simple case where we pick a multicurve
$\lambda \in ML_{\mathbb{Z}}(\tau)$, the
associated measure
is just counting the intersection number
of the curve with a perpendicular segment
near
the train track edge.

The following argument appears to be a
converse construction.
Given a solution to the switch equations,
i.e. an element in $\omega(\tau)$, we can
produce
a multicurve $\lambda \in
ML_{\mathbb{Z}}(\tau)$.
(The method is essentially putting as many
strands
as the weights say near the vertex;
this yields formal sums of curves.)
This also works for formal sums with rational
coefficients.
And for non-negative real coefficients, we
take a thickening of $\tau$.

\begin{theorem}[Thurston]
\label{theorem-thurston-train-tracks}
The weights on $\tau$ uniquely determine a
transverse measure on $L$.
\end{theorem}

\medskip\noindent
Facts:
\begin{enumerate}
\item Every $\lambda \in MS(S)$
is carried by some train track $\tau$ and
determines
a unique point in $\omega(\tau)$.
\item There are finitely many train tracks
$\tau_i$ such that
every $\lambda \in ML(S)$ is carried by at
least $\tau_i$.
From this it follows that we can write
$ML(S)=ML(\tau_1)\cup  ML(\tau_2)\cup
\ldots\cup ML(\tau_n)$
and $ML(\tau)\approx \omega(\tau)\approx
\mathbb{R}^{E-V}$
(under very mild conditions on the train
tracks).

Here is an exercise: suppose $S$ is closed.
Show that $E-V \leq  6g-6$, and
$E-V=6g-6$ if $\tau$ is maximal (which means
that all complementary
regions are triangles).

\item $ML(\tau_i)\cap ML(\tau_j)$ for $i \neq
j$ has dimension
$<6g-6$.
\end{enumerate}

Given those facts,
$$
\begin{aligned}
ML(S)&=ML(\tau_1)\dot\cup \ldots \dot\cup
ML(\tau_n)\\
&=\omega(\tau_1)\dot\cup \ldots\dot\cup
\omega(\tau_n).
\end{aligned}
$$

\noindent
In particular, $ML(S)\approx
\mathcal{E}^{6g+6L+2r}$
($\omega(\tau_1)\approx\mathbb{R}^{6g+6L-2r}
$).
We may now introduce a measure on the space
of laminations:

\begin{definition}
\label{definition-thurston-measure}
The {\it Thurston measure} on $ML(S)$ is
$$
\lim_{L\to\infty}\frac{1}{L^{6g-6+2r}}
\sum_{\lambda
\in ML_{\mathbb{Z}}(S)}
\delta_{\frac{1}{L}\lambda}
$$
\end{definition}

\noindent
One checks that this is a Lebesgue measure
(in the limit) as follows:
$$
\begin{aligned}
ML_{\mathbb{Z}}(S)&=\bigcup_{i=1}
^nML_{\mathbb{Z}}(\tau_i)\\
&\frac{1}{L^{6g-6+2r}}\sum_{\lambda \in
ML_{\mathbb{Z}}\tau_1}
\delta_{\frac{1}{L}\lambda}\\
&\frac{1}{L^{6g-6+2r}}\sum_{(\omega)\in
\omega_{\mathbb{Z}}(\tau_i)}
\delta_{\frac{1}{L}\omega}\\
&\frac{1}{L^{6g-6+2r}}\sum_{p \in
\mathbb{Z}^{6g-6+2r}}
\delta_{\frac{1}{L}p}
\end{aligned}
$$

\noindent
on $\mathbb{R}^{6g-6+2r}$.

Notice that there is an action of the mapping
class group
on the space of laminations: a multicurve is
taken to another multicurve.

The following theorem was used for proving
the speaker's results:

\begin{theorem}[Masur '85]
\label{theorem-masur-85}
$m_{Thu}$ is $\text{Map}(S)$-invariant
and is ergodic with respect to
$\text{Map}(S)$.
\end{theorem}

\medskip\noindent
Now we discuss (geodesic) currents.
Consider de space $\mathcal{G}(\tilde{S})$
of unoriented geodesics on
$\tilde{S}=\mathbb{H}^2$
as $(S^1\times S^1\setminus\Delta)/(a,b)\sim
(b,a)$.
This  is well defined (even topologically,
i.e.
independently of the choice of metric)
since two geodesics on Poincaré plane are
determined by their endpoints.

\begin{definition}
\label{definition-geodesic-current}
A {\it (geodesic) current} is a
$\pi_1(S)$-invariant
Radon measure on $\mathcal{G}(S)$.
\end{definition}

\begin{example}
\label{example-multicurves-are-currents}
Any (multi-)curve is a current.
\end{example}

\noindent
We shall see that not only curves, but also
measured laminations
and in fact the Teichmüller space are
contained in currents:
$$
\begin{aligned}
\{\text{curves}\}&\subset C(S)\\
ML(S)&\subset C(S)\\
\mathcal{T}(S)&\subset C(S)
\end{aligned}
$$

Let $X$ be a hyperbolic metric on $S$.
Let $U=\{\lambda \in
ML(S):\ell_X(\lambda)\leq 1\}$.
If we let $m^L(U_X)=\frac{1}{L^{6g-6}}
\#\{\gamma \in
ML_{\mathbb{Z}}(S)|\ell_X(\gamma)\leq  L\}$,
then $m^L(U_X) \to m_{\text{Thu}}(U_X):=B(X)$
 as $L \to \infty$.

\medskip\noindent
In the following, we think os simple
curves as measured laminations and currents.
Recall that a measured lamination is
a union of disjoint curves.
$$
\xymatrix@C=4em{
ML_{\mathbb{Z}}(S)\ar@{^{(}->}[r]&\overline{
\{\text{simple
multicurves}\}}\ar@{^{(}->}[d]\ar@{^{(}->}[r]
^-{\substack{\text{equality for} \\
\text{closure only!}}}
&ML(S)\ar@{^{(}->}[d]\\
&\overline{\{\text{multi-curves}\}}\ar@{^{(}->}[r]^-{
\substack{\text{equality for} \\
\text{closure only!}}}
&CS=\{\text{geodesic currents}\}.
}
$$

\medskip\noindent
Facts:
\begin{enumerate}
\item (Banahon.) The geometric intersection
number of curves
extends to
$$
i:C(S) \times C(S) \to \mathbb{R}_{\geq 0}
$$
continuously, bi-linear,
$\text{Map}(S)$-invariant.

\item $ML(S)=\{\mu \in
C(S):i(\mu,\mu)=0\}\rightsquigarrow
m_{\text{Thu}}$ measure on $C(S)$
(supported on $ML(S)$).

\item $C_{\text{fin}}(S)=\{\mu \in C(S):
(\mu,\lambda)>0 \forall \lambda \in C(S)\},$
the {\it filling currents}, is an open dense
subset of $C(S)$.

\item (Banahan.) For all $X$ hyperbolic
metric on $S$,
there exists a unique $\lambda_X \in C(S)$
such that
$i(\lambda_X,\gamma)=\ell_X(\gamma)$ for all
curves $\lambda$
(Liouville current) $\rightsquigarrow $
$\mathcal{T}(S) \subset C(S)$.
This way we have a notion of hyperbolic
length of currents $\mu \in C(S)$:
$\ell_X(\mu)=i(\lambda_X,\mu)$.
\end{enumerate}

\begin{definition}
\label{definition-nice-length-function}
A length function on curves on $S$
is {\it nice} if it extends to a function
$F:C(S) \to \mathbb{R}_{\geq 0}$ wich is
continuous, homogeneous
($F(c\cdot\mu)=cF(\mu)\;\forall c>0$,
positive ($F(\mu)>0\;\forall \mu\neq 0$.
\end{definition}

\begin{example}
\label{example-nice-length-functions}
Following are examples of nice length
functions:
\begin{itemize}
\item Hyperbolic length.
\item (Otal.) Negatively curved Riemannian
metric.
\item (E-Parlier-Souto.) (Singular)
Riemannian metric.
\item (Duchen-Leinhein-Rafi.) Euclidean cone
surface.
\item (Banachan.) $\sigma \in C_{fill}(S)$,
$\ell_{\sigma}=(\sigma,M)$.
\item Word length w.r.t. finite generating
set $\pi_1(S)$.
\end{itemize}
\end{example}

\medskip\noindent
Let us go back to our original problem
of counting curves of a specific type.
Fix $\gamma_0$ a (multi-)curve (not
necessarily simple).
$\gamma$ is of type $\gamma_0$ if
$\gamma \in \text{Map}(S)\cdot \gamma_0$.

We want asymptotics of
$\#\{\gamma\text{ type
}\gamma_0:\ell_X(\gamma)\leq L\}$.
Fix $\gamma_0$ and $L>0$. Consider
$m^L_{\gamma_0}=\frac{1}{L^{6g-6}}
\sum_{\gamma\text{ type
}\gamma_0}\delta_{\frac{1}{L}\gamma}$
measure on $C(\gamma)$.

Note that $m^L_{\gamma_0}(U_X)=
\frac{1}{L^{6g-6}}\#
\{\gamma\text{ type
}\gamma_0:\ell_X(\gamma)\leq L\}$.

We are ready to state the main theorem:

\begin{theorem}[Mirzakhani '08, E-Souto '22]
\label{theorem-mirzakhani-e-souto}
For any $\gamma_0$ on $S=S_g$
$$
m^L_{\gamma_0}\to\frac{c(\gamma_0)}{b_g}
m_{\text{Thu}},\qquad
L\to \infty.
$$
\end{theorem}

\noindent
Mirzakhani proved it for simple curves,
while we now presented a general version
using currents.
We write the constants as
$$
\begin{aligned}
c(\gamma_0)&=m_{\text{Thu}}(\{\lambda \in
ML(S)/\text{Stab}(\gamma_0)
|i(\lambda,\gamma_0)\leq 1\},\\
b_g&=\sum_{\alpha_0 \in
ML_\mathbb{Z}/\text{Map}(S)}c(\alpha_0).
\end{aligned}
$$

\noindent
Mirzakhani used another expression.
This one has the virtue of being purely
combinatorial.

Mirzakhani showed the following example.
Consider the genus 2 surface $S_2$.
Then
$$
\lim_{L \to \infty}
\frac{\#\{\text{non-separating
}\gamma|\ell_X(\gamma)\leq L\}}{
\#\{\text{separating
}\gamma|\ell_X(\gamma)\leq L\}}
=48.
$$
This may be interpreted by saying that when
taking a closed simple geodesic,
the probability of it being separating is
$\frac{1}{48}$.

\section{The unstable flow of a partially 
hyperbolic diffeomorphism with an 
application to measure rigidity}
\label{section-unstable-flow-hyp-dyn}

\noindent
Bruno Santiago (UFF?)

Qui 21 mai 2026, 15:30

\medskip\noindent
{\bf Abstract.}
Let $f$ be a partially hyperbolic
diffeomorphism preserving a splitting
$
TM = E^s \oplus E^c \oplus E^u,
$
where the unstable bundle $E^u$ is
one-dimensional.
When $f$ is homogeneous, like the time-one
map of the geodesic flow on a surface of
constant negative curvature or a linear
automorphism of the torus, the unstable
foliation admits a natural parametrization
that makes invariant measures with
absolutely continuous conditionals
(the so-called $u$-Gibbs measures)
correspond exactly to invariant measures
of this flow.
How can this picture be generalized when
$f$ is not homogeneous?
In this talk, I will explain a construction
of a suitable suspension flow of $f$ for
which the unstable foliation becomes the
orbit foliation of a new flow, called the
unstable flow.
Ergodic $u$-Gibbs measures are in
one-to-one correspondence with measures
simultaneously invariant under both flows
and ergodic for the suspension.
As an application, I will deduce the
following measure rigidity statement:
if $E^c$ is expanding and the bundle
$E^s \oplus E^u$ is integrable, then any
$u$-Gibbs measure whose entropy exceeds its
unstable entropy is necessarily an SRB
measure.
This is joint work with Sébastien
Alvarez, Sylvain Crovisier, Martin
Leguil, and Davi Obata.

\begin{remark}[AI-generated overview]
\label{remark-ai-overview-seminar-unstable-flow}
The talk was about translating a geometric
object into a dynamical one. We start with a
partially hyperbolic diffeomorphism
$$
f:M\to M.
$$
This means that there is a $Df$-invariant
splitting
$$
TM=E^s\oplus E^c\oplus E^u.
$$
Here $Df$-invariant means that
$$
Df_x(E_x^s)=E_{f(x)}^s,\qquad
Df_x(E_x^c)=E_{f(x)}^c,\qquad
Df_x(E_x^u)=E_{f(x)}^u.
$$
Thus $f$ moves points, while $Df$ moves the
corresponding tangent directions in a way
compatible with the splitting.

\medskip\noindent
The key simplifying hypothesis is that
$E^u$ is one-dimensional. Hence the unstable
distribution integrates to curves
$$
W^u(x),
$$
called unstable leaves. The first geometric
idea is that these leaves should not merely
be thought of as a foliation, but as the
orbits of a flow. Thus one wants to construct
an unstable flow $\varphi_t^u$ whose orbit
through $x$ is exactly $W^u(x)$.

\medskip\noindent
The nontrivial point is the choice of
parametrization
[because the passage from the orbits
to a vector field is not integration,
but differentiation, which in principle
is easy --- once you have
the right parametrization!]. 
A one-dimensional leaf can
be parametrized in many different ways, and
the dynamics of $f$ selects a preferred one.
If $x$ and $y$ lie on the same unstable leaf,
the function
$$
\rho_x^u(y)
$$
compares the expansion of $Df^n$ in the
unstable direction at $x$ and at $y$. Thus
$\rho_x^u$ is a dynamical density along
$W^u(x)$.

\medskip\noindent
This density defines a coordinate on the
unstable leaf:
$$
H_x^u(y)
=
\int_x^y
\rho_x^u(z)\,d\operatorname{Leb}_x^u(z).
$$
The inverse map
$$
\Phi_x^u=(H_x^u)^{-1}
$$
then gives the natural parametrization of the
unstable leaf. This is the geometric heart of
the construction.

\medskip\noindent
The measure-theoretic part enters because in
dynamics one studies not only where an orbit
goes, but also its long-time statistical
distribution. For a point $x$, the empirical
measures
$$
\mu_{n,x}
=
\frac{1}{n}
\sum_{j=0}^{n-1}
\delta_{f^j(x)}
$$
record the proportion of time the orbit of
$x$ spends in different regions of $M$.
Equivalently, for a measurable subset
$A\subset M$, one has
$$
\mu_{n,x}(A)
=
\frac{1}{n}
\#\{0\leq j<n:f^j(x)\in A\}.
$$
Thus $\mu_{n,x}(A)$ is the proportion of
time, up to time $n$, that the orbit of $x$
spends inside $A$. This is where time enters:
we are not measuring the size of $A$ itself,
but rather how frequently the forward orbit
$$
x,\ f(x),\ f^2(x),\ldots
$$
visits $A$.

\medskip\noindent
A $u$-Gibbs measure is an invariant measure
which is smooth in the unstable direction.
More precisely, if
$$
\{\mu_x^u\}
$$
is the disintegration of $\mu$ along unstable
leaves, then
$$
\mu_x^u\ll \operatorname{Leb}_x^u.
$$
So the measure may be globally complicated,
but leaf-by-leaf along $W^u$ it has a density
with respect to ordinary leafwise length.

\medskip\noindent
The main theorem of the talk says that such
$u$-Gibbs measures can be characterized after
passing to a suspension space. The original
map $f$ is discrete-time, while the suspension
turns it into a continuous-time flow
$\psi_t$. The theorem identifies $u$-Gibbs
measures with lifted measures which are
invariant both under the suspension flow
$\psi_t$ and under the unstable flow
$\varphi_t^u$.

\medskip\noindent
Finally, the application concerns SRB
measures. An SRB measure is stronger than a
$u$-Gibbs measure: it is meant to describe
the statistics seen by a positive-Lebesgue
set of initial conditions. In the theorem
from the end of the talk, if the center
direction is expanding and a $u$-Gibbs measure
has more entropy than what is accounted for
only by the unstable direction, namely
$$
h_\mu(f)>
\int_M \log \lambda_x^u\,d\mu(x),
$$
then the measure must be SRB.

\medskip\noindent
In the language of the Heap Project, this is
not like constructing a scheme from equations;
it is more like extracting an invariant
object from asymptotic data. Just as a Hilbert
polynomial records asymptotic information
about graded pieces, an invariant measure
records asymptotic information about the
distribution of an orbit. The slogan of the
talk is therefore:
$$
\text{unstable expansion}
\quad\leadsto\quad
\text{dynamical density}
\quad\leadsto\quad
\text{unstable flow}
\quad\leadsto\quad
\text{measure rigidity}.
$$
\end{remark}

\noindent
Let $S:M \to M$ be a $C^r(M)$ map.
Let $V$ be a vector field
inducing a foliation
whose leaves form a $C^r$ foliation
on $M$.
The idea is
that this dimension-1 
foliation is dilated by the
diffeomorphism.

\begin{example}
\label{example-sl3-spectrum}
\begin{itemize}
\item $A \in \text{SL}_3(\mathbb{Z})$,
with spectrum
$\text{spec}(A)
= \{0<\lambda_A^s<1<
\lambda^c_A
<\lambda^u_A\}$.
$W_A^u
= $ folitation with orbits
of irrational flux $\varphi_t$
of $f_A:\mathbb{T}^3 \to \mathbb{T}^3$.
\end{itemize}
\end{example}

\begin{definition}[AI]
\label{definition-u-gibbs-measure}
Let $f:M \to M$ be a partially hyperbolic
diffeomorphism, and let
$\mathcal F^u$ be its unstable foliation.
A measure
$$
\mu \in \mathcal P_f^{\operatorname{reg}}(u)
$$
is called a $u$-Gibbs measure if its
conditional measures along unstable leaves
are absolutely continuous with respect to
Lebesgue measure on the leaves.

More precisely, if
$$
\{\mu_x^u\}_{x \in M}
$$
denotes the disintegration of $\mu$ along
the unstable foliation, then
$$
\mu_x^u \ll \operatorname{Leb}_x^u
$$
for $\mu$-almost every $x \in M$.
\end{definition}

\begin{remark}
\label{remark-stable-directions-dyn}
$$
TM = E^s \oplus E^c \oplus E^u.
$$
\end{remark}

\begin{remark}
\label{remark-leb-ae-ai}
For Lebesgue almost everywhere,
$x \in M$, $n_k \to \infty$,
$$
\mu_{n_k,x}
:=
1 \sum_{n_k=0}^{n_k-1}
\delta_{f_k} \to \mu
$$
which implies that $\mu \in 
\text{Leb}^u$.
\end{remark}

\noindent
I'm very lost, here's AI-corrected
above remark:

\begin{remark}
\label{remark-lebesgue-ae}
The phrase ``Lebesgue almost every $x$''
means almost every point with respect to
the Riemannian volume measure on $M$.

Given $x \in M$, its empirical measures are
$$
\mu_{n,x}
:=
\frac{1}{n}
\sum_{j=0}^{n-1}
\delta_{f^j(x)}.
$$
A measure $\mu$ is called physical if its
basin
$$
B(\mu)
:=
\left\{
x \in M :
\mu_{n,x} \to \mu
\right\}
$$
has positive Lebesgue measure, where the
convergence is weak-$*$ convergence of
probability measures.
\end{remark}

\begin{lemma}[Ledrappier density]
\label{lemma-ledrappier}
$$
\frac{d\mu^u_x}{
d\text{Leb}_x^u}
=c(X)\rho_x^u
$$
where $\rho_x^u:w^u(x)\to \mathbb{R}$,
$$
\rho_x^u(y)
=
\lim_{n \to \infty} 
\frac{\|D\bar{f}^n(y)V^u\|}{
D\bar{f}^n(x) V^u}
$$
is the dynamic density function.
\end{lemma}
\begin{lemma}[Ledrappier density, AI]
\label{lemma-ledrappier-density}
Let $\mu$ be a $u$-Gibbs measure. Then the
conditional measure $\mu_x^u$ is absolutely
continuous with respect to $\operatorname{Leb}_x^u$,
and its density is of the form
$$
\frac{d\mu_x^u}{d\operatorname{Leb}_x^u}(y)
=
c(x)\rho_x^u(y),
$$
where $c(x)>0$ is a normalizing constant and
$$
\rho_x^u:W^u(x)\to \mathbb R
$$
is given by
$$
\rho_x^u(y)
=
\lim_{n\to\infty}
\frac{
\|D\bar f^{-n}(y)|_{E_y^u}\|
}{
\|D\bar f^{-n}(x)|_{E_x^u}\|
}.
$$
Equivalently, in the one-dimensional unstable
case,
$$
\rho_x^u(y)
=
\lim_{n\to\infty}
\frac{
\|D\bar f^n(x)v_x^u\|
}{
\|D\bar f^n(y)v_y^u\|
},
$$
where $v_x^u\in E_x^u$ and $v_y^u\in E_y^u$
are unit vectors.
\end{lemma}

\medskip\noindent
{\bf Properties of EDD.}
\begin{itemize}
\item 
$$
\rho_y^(u)
=\frac{j}{\rho^u_x(y)}
$$
\item 
$$
z,y \in W^u(x)
\implies 
\rho^u_x(z)
=\rho_x^u(y)\rho_y^u(z)
$$
\item 
$$
\rho^u_{f(x)}(f(y))
=\frac{\lambda_x^u}{\lambda_y^u}
\rho_x^u(y).
$$
\end{itemize}

\medskip\noindent
{\bf Affine structure.} 
$$
H_x^u:W^u(x) \to \mathbb{R}
$$
$$
H_x^u(y)
=\int_x^y \rho_x^u(z)
d\text{Leb}_x^u(z)
$$
$$
\Phi_x^\mu
:=(H_x^u)^{-1}
:\mathbb{R} \to W^u(x).
$$

Change of coordinates:
$$
H_y^u \circ\Phi_x^u(t)
=\rho_y^u(x)H
+H_y^u(x).
$$
\begin{definition}[Unstable flux]
\label{definition-unstable-flux}
$$
\begin{aligned}
\tilde{\varphi}_t:
M \times \mathbb{R} &\longrightarrow 
M \times \mathbb{R} \\
(x,\alpha) &\longmapsto 
(\Phi^u_x(e^{-x}t),
\alpha
+\text{log}\rho_x^u\Phi_t^u(e^{-\alpha}t))
\end{aligned}
$$
\end{definition}

\noindent
Let's check why this, in fact,
defines a flux:

\begin{lemma}
\label{lemma-defines-a-flux}
$$
\tilde{\varphi}_t \circ \tilde{\varphi}_s
(x,\alpha)=\tilde{\varphi}_{t+s}(x,\alpha)
$$
\end{lemma}

\begin{proof}
Provided
\end{proof}

\begin{theorem}
\label{theorem-gibbs-measure}
$\mu \in \text{Gibbs}_{\text{reg}}^u(f)$,
then there exists a unique
$$
\hat{\mu} 
\in \mathcal{P}_{\psi}^{\text{reg}}
(\hat{\mu}) \cap
\mathcal{P}_\varphi(\hat{\mu})
$$
\end{theorem}

$$
D(f)
=\{(x,\alpha):0\leq \alpha<\lambda_x^u\}
\subset\hat{\mu}
$$
$$
\hat{\mu}
=\frac{1}{
\int_M \text{log} \lambda_x^u
d\mu(x)}\mu\otimes\text{Leb}_{\mathbb{R}}
(\mathbb{T}^{-x}(A)).
$$

\medskip\noindent
{\bf Applications of Theorem
\ref{theorem-gibbs-measure}}

\noindent
Back to the totally partially hyperbolic
context, suppose there is a decomposition
$$
TM^3 = E^s \oplus E^c \oplus E^u
$$
such that
$$
0 \leq \|Df(x)|_{E^s}\|
<I<\|Df(x)|_{E^c}\|
<\|Df(x)|_{E^u}\|.
$$
Then $f$ is Anosov partially hyperbolic.
Thus, there exist foliations
$W^s,W^c,W^u$ such that
$E^s,E^c$ and $W^u$ and
$W^{cs}, W^{du}$
$$
F^s \oplus E^c \oplus E^c \oplus E^u.
$$

\begin{definition}
\label{definition-}
$\mu \in \mathcal{P}_f(u)$
is SRB if all the conditionals
$$
\{\mu_x^{c\mu}\}_{x \in M}
$$
along $W^{c\mu}(\alpha)$
satisfy
$$
\mu_x^{c\mu}\ll \text{Leb}_x^{c\mu}.
$$
\end{definition}

\begin{theorem}[Sinoi-Rulh-Bouso]
\label{theorem-srb}
There exists a unique SRB measure
$\mu$.
\end{theorem}

\noindent
Conjecture (ALOS, Katz): 
if $E^s \oplus E^u$
are not integrable then any
$$
\mu \in \text{Gibbs}_{\text{reg}}^{\mu}(f)
$$
is SRB.

\begin{theorem}[B]
\label{theorem-gibbs-b}
If $E^s \oplus E^u$ is integrable
and
$$
\mu \in \text{Gibbs}_{\text{reg}}^u(u)
$$
satisfies
$$
h_\mu(f)
> \int \text{log}\lambda_x^u d\mu(x),
$$
then $\mu$ is SRB.
\end{theorem}

\medskip\noindent
{\bf Some ideas.}

 \noindent
$x \in M \mapsto \text{log}\lambda_x^u$
is $C^{1+\alpha}$.

Descição das $W^*$
$$
\widetilde{W}^s(x,\alpha)
=\{(x,\alpha+\theta_x^{s}(y)
:y \in W^s(x)\}
$$
where
$\theta(\alpha)=\text{log}\lambda_x^u$
and
$\theta_x^s(y)
=\sum_{\ell=0}^{\infty}
\theta(f^\ell y)
-\theta(f^\ell x)$.

\section{Weil--Petersson homeomorphisms 
and maximal surfaces in anti-de Sitter space}
\label{section-weil-petersson}

\noindent
Graham Smith

May 21, 2026.

\medskip\noindent
{\bf Abstract.}
A classe de Weil--Petersson é uma classe
notável de homeomorfismos do círculo que se
encontra na fronteira de vários campos da
matemática contemporânea, incluindo a teoria
de Teichm\"uller, a teoria de
Schramm--Loewner, a geometria hiperbólica e
assim por diante.
Neste trabalho, apresentamos uma nova
perspectiva, estudando essa classe do ponto
de vista da geometria anti-de Sitter.
Mostramos que um homeomorfismo do círculo é
Weil--Petersson se e somente se seu gráfico
limitar uma superfície máxima completa,
semelhante ao espaço, em
$\mathrm{AdS}^{2,1}$ de área finita
renormalizada.

Este é um trabalho conjunto com F. Diaf,
A. Moriani, R. Smai e E. Trebeschi.
2604.17804.

\medskip\noindent
We are interested in Weil-Petersson
Teichmüller space $\mathcal{T}_0$
This is a very interesting object
that intersects several branches of 
mathematics:
\begin{itemize}
\item Renormalized area of minimal surfaces
in $\mathbb{H}^3$.
\item Stochastic Loewner evolution.
\item Geometric computer vision.
\item Berezin Quantization and Bott-Virasoro
group.
\item Planar Coulomb gasses.
\end{itemize}

\medskip\noindent
{\bf Loewner transport.}
Let  $\Omega \subset \hat{\mathbb{C}}$
be a Jordan domain (interior of Jordan
curve).
Let $p,q \in \Omega$.
For example, $\Omega=\mathbb{H}^2$,
$p=0$ and $q=\infty$.
Let $\gamma:[0,1]\hookrightarrow
\overline{\mathbb{H}^+}$
with $\gamma(0)=0$ and $\gamma(t) \in
\mathbb{H}^+$ for all $t>0$.

For all $t$ define
$$
\Omega_t:\mathbb{H}^+\setminus
\gamma([0,t])
$$
$$
g_t:\mathbb{H}^t \to \Omega_t
$$
be the unique uniform map
$$
g_t(z)
=z + \frac{2\varphi(t)}{z}
+ o \left(\frac{1}{z}\right)
$$
where $\varphi$ is real valued and
increasing. Parametrize $\gamma$
by $\varphi$:
$$
g_t(z)=z+ \frac{2t}{z}
+o \left(\frac{1}{z}\right).
$$
Caratheory:
$$
W(t)=g^{-1}_t(\gamma(t))
$$
$$
W:\underbrace{[0,L]}_{\text{domain of }
\gamma} \to \mathbb{R}.
$$

The {\it Lowener transform} (1923)
is the map
$$
\gamma \mapsto W
$$
and $W$ is the {\it driving function},
invertible.

\begin{remark}
\label{remark-brownian-loewner}
If $W$ is Brownian we obtain a 
stochastic loewner evolution.
Yilin Wane.
\end{remark}

\medskip\noindent
{\it Loewner energy:}
$$
I(\gamma)
:= \frac{1}{2}
\int_0^L W'(t)^2 dt
$$
\begin{remark}
\label{remark-loewner-energy-independent}
Jordan curves have a well-defined
Loewner enery
[some remark I didn't understand concerning
$\Gamma$ Jordan curve, $p,q \in \Gamma$…]
\end{remark}

\noindent
In this talk we shall take this
as a definition:

\begin{theorem}[Wang]
\label{theorem-wang}
If $\Gamma \subseteq \hat{ \mathbb{C}}$
is a Jordan curve then
$I(\Gamma) < \infty$ if and only if
$\Gamma$ is in $\mathcal{T}_0$,
the {\it Weil-Petersson Teichmüller space}.
\end{theorem}

\medskip\noindent
{\bf Universal Teichmüller space.}
Let
$B_1^{\infty}$ be the open unit ball in
$L^\infty(\mathbb{H}^+)$.
Let $\mu \in B_1^\infty$,
$$
\mu_0(z)=
\begin{cases}
\mu(z)\qquad &\text{ in }\mathbb{H}^+ \\
0\qquad &\text{ in }\mathbb{H}^-.
\end{cases}
$$

The upshot of the following theorem
is that these functions behave
like holomorphic functions:

\begin{theorem}[Measurable Riemann Mapping
Theorem]
\label{theorem-measurable-rmt}
There exists a unique $f_\mu:\hat{\mathbb{C}}
\to \hat{\mathbb{C}}$ such that
\begin{enumerate}
\item $f_\mu(0)$, $f_\mu(1)=1$,
 $f_\mu(\infty)=\infty$.
\item 
$$
\frac{f_{\mu,\bar{z}}}{
f_{\mu z}}=\mu_0.
$$
\end{enumerate}
\end{theorem}

\noindent
Recall:

\begin{definition}
\label{definition-quasiconformal-map}
$f:\Omega(\subseteq\mathbb{C})
\to \mathbb{C}$
is {\it quasiconformal} if
\begin{enumerate}
\item $f$ is a local homeomorphism,
\item $f$ has locally $L^2$ 
dist. derivatives.
\item There exists an $L^\infty$
function $\mu$ such that
$\|\mu\|_{L^\infty}<1$ such that
$$
f_{\bar{z}}=f_z\cdot \mu
$$
and $\mu$ is the {\it Beltrami differential} 
of $f$.
\end{enumerate}
\end{definition}

\noindent
$\mu_0$ is the ``extension by $0$''.
There's also:
$$
\mu_s(z)
=\begin{cases}
\mu(z)\qquad &z \in \mathbb{H}^+ \\
\overline{\mu(\bar{z})}\qquad &
z \in \mathbb{H}^-
\end{cases}
$$
and by the Measurable Riemann Mapping 
Theorem, there exists a unique
$f^u:\hat{\mathbb{C}} \to \hat{\mathbb{C}}$
such that $\sim$
$$
f_{\bar{z}}^\mu
=\mu_s f^\mu.
$$

[Picture which I think
says $f_\mu$ maps the upper half plane
to a Jordan region]

\begin{definition}
\label{definition-quasicircle}
A Jordan curve $\Gamma$ is a
{\it quasicircle} whenever
$$
\Gamma=f_\mu(\hat{\mathbb{R}})
$$
\end{definition}

\noindent
Symmetry:
[Picture showing $f^\mu$
putting the upper half plane
to the whole plane.]

$f^\mu$ restricts to a homeomorphism
of $\hat{\mathbb{R}}$,
$$
\varphi:\hat{\mathbb{R}}
\to \hat{\mathbb{R}}
$$
is quasisymmetic whenever
$$
\varphi=f^\mu|_{\hat{\mathbb{R}}}
$$
for some $\mu$.

$$
\xymatrix{
& B_1\subseteq L^\infty
(\mathbb{H}^+)\ar[dl]\ar[dr]\\
\text{Quasicircles}
\ar[rr]_{\text{Homeo}}
&&
\text{Quasisymmetries}
}
$$
where the labelled homeomorpihsm
is the {\it conformal Welding}.
If $\mu_1, \mu_2 \in B_1$, then
$$
\text{Im}(f_{\mu_1})
= \text{Im}(f_{\mu_2})
\iff
f^{\mu_1}|_{\hat{\mathbb{R}}}
=f^{\mu_2}|_{\hat{\mathbb{R}}}
$$

Bers was very interested in this triangle,
and introduced what he called
``universal Teichmüller space''.

$$
\begin{aligned}
\mathcal{T}&=\text{QCirc}/
\text{PSL}(2,\mathbb{C})\\
&=\text{QSymm}/\text{PSL}(2,\mathbb{R})
\times \text{PSL}(2,\mathbb{R})\\
&=B^\infty_1/\sim.
\end{aligned}
$$

Bers showed that
\begin{itemize}
\item $\mathcal{T}$ is a Banach
manifold.
\item Every $TM$ space embeds smoothly
in $\mathcal{T}$.
\end{itemize}

\noindent
Every Teichmüller space has a 
Weil-Petersson metric. Does $\mathcal{T}$
carry a metric which is induced
by a Weil-Petersson metric? No.

Takhtajan-Teo ($\sim$2004):
$\mathcal{T}$ carries a Hilbert
manifold structure $\mathcal{T}_1$,
it has $\mathbb{R}$ connected components,
contains every Teichmüller space
and its metric (induced by the
Hilbert manifold structure)
induces a Weil-Petersson
metric.

\begin{definition}
\label{definition-wp-teichmuller}
$\mathcal{T}_0 \subseteq T_1$
is the connected component of $I_0$.
\end{definition}

\medskip\noindent
{\bf Properties.}

\begin{itemize}
\item $\mathcal{T}_0$ has a
unique homogeneous Kähler metric.
\item Kähler potential is $I$.
\item Among the dozens
of characterizations of this space,
we have:
$\Gamma$ is $\mathcal{T}$
if and only if $\Gamma$
is chord arch and arc-length
parametrization is $H^{3/2}$,
there exists $c$ such that
$$
d_{\text{ext}}(p,q)
\leq  d_{\text{int}}(p,q)
\leq cd_{\text{ext}}(p,q).
$$
\end{itemize}

\noindent
Here's another characterization 
of this space:

\begin{theorem}[Bishop]
\label{theorem-bishop-teichmmuller}
$\Phi \in \mathcal{T}_0$ if and only if
$\Gamma$ bounds a minimal surface
$\Sigma$ in $\mathbb{H}^3$
of finite topology and finite
renormalized area, which
is the integral of the determinant
of the second fundamental form
(and can be infinite).
\end{theorem}

\medskip\noindent
{\bf Reminder on anti de Sitter space.}

$$
\text{AdS}^{2,1}
=
\{(x_1,…,x_4)
: x_1^2+x_2^2-x_3^2-x_4^2=-1\}
/\{\pm  1\}.
$$
Also, define for any $M$
$$
\text{Det}(M):=
-\frac{1}{2}
\text{Tr}(J_0M J_0M^t)
$$
for
$$
J_0=
\begin{pmatrix}0&-1\\ 
1&0\end{pmatrix}.
$$
Then
$$
\begin{aligned}
\text{AdS}^{2,1}&=
\{M : \text{Det}M=0, M\neq 0\}
/\{\pm I_0\}
\end{aligned}
$$

\noindent
and
$$
\begin{aligned}
\delta_{\infty}\text{AdS}^{2,1}&=
\{M: \text{Det}(M)=0,M \neq 0\}\\
&=\{M:\text{rk}M=1\}/\mathbb{R}^*\\
&=\underbrace{\mathbb{R}P^1}_{\text{image}} 
\times \underbrace{\mathbb{R}P^1}
_{\text{kernel}}.
\end{aligned}
$$

\noindent
A homeomorphism 
$\varphi:\mathbb{R}P^1 \to \mathbb{R}P^1$
gives a graph $\Lambda_\varphi
\subset \partial_{\infty}\text{AdS}
^{2,1}$ which is spacelike.

Anti de Sitter space is a lot nicer
than hyperbolic space:

\begin{theorem}[Labourie-Toul-Wolf]
\label{theorem-lab-tou-wol}
Every spacelike curve $\Lambda$
in $\partial_{\infty}\text{AdS}^{2,1}$
bounds a unique complete maximal disk.
\end{theorem}

\begin{theorem}[Graham-et al.]
\label{theorem-graham-teich}
$\varphi$ is WP if and only if
$D_\varphi$ has finite renormalized area.
\end{theorem}

\section{On multiplicities in length spectra of semi-arithmetic surfaces}
\label{section-multiplicities-length-spectra-semi-arithmetic-surfaces}

\noindent
Misha Belolipetsky.

Ter 26 mai 2026, 15:30.

SALA 236.

\medskip\noindent 
{\bf Abstract.} Mostramos
que determinadas superfícies semi-aritméticas
de Riemann têm um crescimento exponencial das
multiplicidades médias em seu espectro de
comprimento. Anteriormente, o crescimento
exponencial das multiplicidades médias era
conhecido apenas para superfícies
aritméticas, embora os resultados
experimentais que indicavam a possibilidade
de exemplos não aritméticos já estivessem
disponíveis. Explicarei alguns detalhes
interessantes de nossa prova e indicarei
algumas questões em aberto relacionadas. A
palestra é baseada em um trabalho conjunto
recente com Gregory Cosac, Cayo Dória e
Gisele Teixeira Paula.

\medskip\noindent
Zoll metrics $\leftrightarrow$
systole has large mutiplicity.

\medskip\noindent
Our case $\leftrightarrow$
``large geodesics have
very large multiplicity''.

\medskip\noindent
Plan of the talk:
\begin{enumerate}
\item Introduction to semiarithmetic
surfaces.
\item 
\end{enumerate}


\medskip\noindent
Let $S=\mathbb{H}^2/\Gamma$
be a Riemann surface, i.e.
the hyperbolic plane modulo a 
Fuschian group, which is a finite discrete
subgroup of $\text{PSL}(2,\mathbb{R})$
of finite coarea.
We introduce the number $\mathcal{N}(\ell)$,
the number of closed geodesics on $S$
of length $\leq \ell$.

\begin{theorem}[Huber, Selberg]
\label{theorem-huber-selberg}
$\mathcal{N}(\ell)
\sim \frac{\text{exp}\ell}{\ell}$.
\end{theorem}

\noindent
This number is well-defined: i.e.
there are only finitely many closed
geodesics in a hyperbolic surface.

\begin{remark}
\label{remark-nlog}
$\mathcal{N}(\text{log} x)
\sim \frac{x}{\text{log} x}$.
\end{remark}

\noindent
The length set is
$$
\{0<\ell_1<\ell_2<…\}
$$
Also we denote
$$
\mathcal{N}'(\ell)
=\#\{n : \ell_n <\ell\}
$$
and $g(\ell)$ is the multiplicity of
$\ell_n$.

\begin{definition}
\label{definition-mean-multiplicity}
The {\it mean multiplicity} 
$\langle g(\ell)\rangle$ is a continuous
monotone function
$$
\mathcal{N}(\ell)
=\sum_{\ell_n<\ell}g(\ell_n)
\sim
\sum_{\ell_n \leq  \ell}
\langle f(\ell)\rangle
$$
\end{definition}

$$
\langle g(\ell)\rangle
\sim
\frac{\mathcal{N}(\ell)}{\mathcal{N}'
(\ell)}.
$$
\medskip\noindent
In a hyperbolic surface
$S = \mathbb{H}^2/\Gamma$,
closed geodesics correspond to
hyperbolic elements $\gamma \in \Gamma$,
$|\text{tr}\gamma|>2$.
Similarly, lengths correspond to traces.
We have
$$
\ell(\gamma)
=2 \text{arcosh}
\left(\frac{|\text{tr}\gamma|}{2}\right).
$$
$$
\mathcal{L}(\Gamma)
=\{|\text{tr}(\gamma)|:\gamma \in \Gamma,
\gamma\text{ hyp.}\}
$$
\begin{example}
\label{example-modular-group}
$\Gamma = \text{PSL}_2(\mathbb{Z})$,
$\mathcal{L}(\Gamma) \subset \mathbb{N}$
(in particular, discrete in $\mathbb{R}$).
\end{example}

$$
\#(\mathcal{L}(\Gamma)
\cap[2,T])<T.
$$
Each length $\ell$
produces a unique trace
$\leq  2 \text{cosh}\left(\frac{\ell}{2}
\right)$.
$$
\mathcal{N}'(\ell)
\leq 2 \text{cosh}(\ell/2)\leq 2e^{\ell/2},
$$
$$
\langle g(\ell)\rangle
\sim\frac{\mathcal{N}(\ell)}{
\mathcal{N}'(\ell)}
\geq \frac{\text{exp}(\ell/2)}{2\ell}
$$

\begin{theorem}[Luo-Sernak (1994)]
\label{theorem-luo-sarnak}
Instead of discretness
of $\mathcal{L}(\Gamma)$ in $\mathbb{R}$,
it is enough
$$
\# (\mathcal{L}(\Gamma)
\cap [N-1,N]
\leq \text{const.}
$$
for any $N \in \mathbb{N}$.
\end{theorem}

\noindent
This is called the
{\it bounded clustering property (BC)}.

\begin{example}
\label{example-triangle24oo}
$\Gamma=\Delta(2,4,\infty)$.
It is well-known this group is generated
by
$$
\begin{pmatrix}0&1\\ -1&0\end{pmatrix},
\qquad 
\begin{pmatrix}1&\sqrt{2}\\ 0&1\end{pmatrix}.
$$
$\mathcal{L}(\Gamma) \supset 2 \mathbb{N}
\cup \sqrt{2}\mathbb{N}$.
Non discrete in $\mathbb{R}$,
but satisfies B. C.
\end{example}

\begin{remark}
\label{remark-bc-holds}
BC holds for ``some''
arithmetic groups.
\end{remark}

\begin{itemize}
\item BC $\implies $ exponential group
of $\langle g(\ell)\rangle$.
\item All arithmetic Fuchsian groups have 
BC.
\item Conjecture: BC iff group is $\Gamma$
is arithmetic.
\item Non-compact case: Gerinska-Leuzinger
(2008).
\item Open in general.
\end{itemize}

\medskip\noindent
Let's try to define arithmetic 
and semi-arithmatic groups.

Let $\Gamma \subset \text{SL}_2(\mathbb{R})$
be a Fuschian group.
Consider the {\it invariant trace field}.
It's the field generated by the traces
of the elements in the group:
$$
K:=\mathbb{Q}\bigl(\{\text{tr}(\gamma^2):
\gamma \in \Gamma\}\bigr).
$$

Recall that a totally real field
is such that all complex embeddings are real.
More precisely:

\begin{definition}
\label{definition-totally-real-field}
A number field $K$ is {\it totally real} if 
every field embedding
$K \hookrightarrow \mathbb{C}$ has image 
contained in $\mathbb{R}$.
\end{definition}

\begin{definition}
\label{definition-semiarithmetic}
Let $\Gamma$ be a Fuchsian group
of finite covolume. $\Gamma$ is called
{\it semiarithmetic} is $K$ is totally
real and
$\text{tr}(\gamma) \in \mathcal{O}_K$
--- the integers of $K$.
\end{definition}

$$
\Gamma^{(2)}
\subset \text{SL}_2(\mathcal{O}_K)
\underbrace{\subset}_{\text{lattice}}
\text{SL}_2(K \otimes_\mathbb{Q}\mathbb{R})
=\text{SL}_2(\mathbb{R})^2
\times 
\underbrace{\text{SO}(3)^{[K:\mathbb{Q}]-r}}_{
\text{compact}}
$$
That is $\text{SL}_2(\mathbb{R})^r$
is a noncompact semisimple piece,
and $\text{SO}(3)^{[K:\mathbb{Q}]-r}$
is compact. Then we take projection,
which ammounts to ``modding out the
noncompact part''.

  $$
  \xymatrix{
  \mathrm{SL}_2(K \otimes_{\mathbb{Q}}
  \mathbb{R}) \ar[r]^-{\pi_{\infty_1}}
  \ar[dr]_{\pi_{\infty_2}} &
  \mathrm{SL}_2(\mathbb{R})^r
\times\text{SO}(3)^{[K:\mathbb{Q}]-r} \\
  & \mathrm{SL}_2(\mathbb{R})
  }
  $$

Now for $r = 1$ corresponds with 
$\Gamma$ being arithmetic.
For  $r>1$, is semi-arithmetic,
$r$ is the arithmetic dimension.
More precisely,

\begin{definition}
\label{definition-arithmetic-group}
Let $\Gamma$ be a Fuchsian group of finite
covolume, and let $K$ be its invariant trace
field. We say that $\Gamma$ is
{\it arithmetic} if $\Gamma^{(2)}$ is
commensurable with $\mathrm{SL}_2(\mathcal{O}_K)$.
\end{definition}

What is happening here is that one starts
with a group of matrices over the number
field $K$ and then views it inside the much
larger real Lie group
$\mathrm{SL}_2(K \otimes_\mathbb{Q}\mathbb{R})$.
After extending scalars, this ambient group
splits into one or more noncompact
$\mathrm{SL}_2(\mathbb{R})$ factors and a
product of compact factors. The projection
diagram is the way of choosing the real
factor that gives the actual action on the
hyperbolic plane, while the other factors are
auxiliary arithmetic data. In the case
$r=1$, there is only one noncompact factor,
and that is exactly the arithmetic case;
when $r>1$, the same construction gives the
semi-arithmetic situation.

\medskip\noindent
Now,
$$
\xymatrix{
\{|\text{tr}\gamma|:\gamma \in \Gamma^{(2)}\}
\ar[r]^{\sigma_i}\ar[dr]&  \text{unbounded}\\
&C[0,2]
}
$$
satisfies BC by the pigeon hole.

\begin{example}[Golden triangle]
\label{example-golden-triangle}
$$
\Gamma = \Delta(2,5,\infty)
=\langle \begin{pmatrix}0&-1\\
-1&0\end{pmatrix},
\begin{pmatrix}1&\gamma\\ 0&1\end{pmatrix}
\rangle,
$$
where $\gamma=\frac{1+\sqrt{5}}{2}$.
Then
$$
\Gamma\underbrace{\subset}_{\text{thin}}
\text{SL}_2(\mathbb{Z}[\sqrt{5}])
\underbrace{\subset}_{\text{arithm.}}
\text{SL}_2(\mathbb{R})^2.
$$
Here $K=\mathbb{Q}(\sqrt{5})$,
$\sigma_1=\text{id}:K \hookrightarrow 
\mathbb{R}$,
$\sigma_2:\sqrt{5} \to -\sqrt{5}$.
$\sigma_1(\Gamma) = \Gamma
\subset\text{SL}_2(\mathbb{R})$ - Fuschian
subgroup.
$\sigma_2(\Gamma)\subset
\text{SL}_2(\mathbb{R})$ - not discrete.
\end{example}

\noindent
And that's it! Since this was particularly
difficult to grasp, let us
include an AI-summary:

\medskip\noindent
{\bf On multiplicities in length spectra of
semiarithmetic surfaces.}
This was a talk by Misha Belolipetsky on
joint work with Gregory Cosac, Cayo Dória
and Gisele Teixeira Paula. The main result
is that certain semiarithmetic Riemann
surfaces have exponential growth of the
mean multiplicities in their length spectra.
Previously this phenomenon was known for
arithmetic surfaces; the point of the talk
was that it also occurs in some genuinely
semiarithmetic examples.

\medskip\noindent
The geometric object is a hyperbolic surface
$$
S=\mathbb{H}^2/\Gamma,
$$
where $\Gamma$ is a Fuchsian group, i.e. a
discrete subgroup of
$\operatorname{PSL}_2(\mathbb{R})$ of finite
coarea. We write $\mathcal{N}(\ell)$ for the
number of closed geodesics on $S$ of length
at most $\ell$. The Huber--Selberg theorem
says
$$
\mathcal{N}(\ell)\sim \frac{e^\ell}{\ell}.
$$
Thus the total number of closed geodesics
grows exponentially. The question of the talk
is subtler: if we list only the distinct
lengths
$$
0<\ell_1<\ell_2<\cdots,
$$
and write
$$
\mathcal{N}'(\ell)=\#\{n:\ell_n<\ell\},
$$
then how large is the average multiplicity
of a length? If $g(\ell_n)$ denotes the
number of closed geodesics with length
$\ell_n$, then morally
$$
\langle g(\ell)\rangle
\sim
\frac{\mathcal{N}(\ell)}{\mathcal{N}'(\ell)}.
$$
So the problem becomes: there are very many
closed geodesics, but how many different
lengths do they produce?

\medskip\noindent
The bridge from geometry to arithmetic is
given by traces. Closed geodesics correspond
to hyperbolic conjugacy classes in
$\Gamma$. If $\gamma\in\Gamma$ is
hyperbolic, then $|\operatorname{tr}\gamma|>2$
and
$$
\ell(\gamma)
=
2\operatorname{arccosh}
\left(\frac{|\operatorname{tr}\gamma|}{2}\right).
$$
Thus length is controlled by trace. In
particular, many different conjugacy classes
can have the same trace, hence the same
length. For example, in the modular group
$\operatorname{PSL}_2(\mathbb{Z})$, traces
are integers. Therefore the number of
possible trace values up to size
$2\cosh(\ell/2)$ is at most on the order of
$e^{\ell/2}$. Combining this with
Huber--Selberg gives
$$
\langle g(\ell)\rangle
\geq
\frac{e^\ell/\ell}{e^{\ell/2}}
=
\frac{e^{\ell/2}}{\ell}.
$$
This is the basic mechanism: if arithmetic
forces the trace set to be sparse, then many
geodesics must share the same length.

\medskip\noindent
Luo and Sarnak observed that one does not
need the trace set to be literally discrete
in $\mathbb{R}$. It is enough to have the
bounded clustering property:
$$
\#(\mathcal{L}(\Gamma)\cap[N-1,N])
\leq C
$$
for all integers $N$. Here
$\mathcal{L}(\Gamma)$ is the set of absolute
trace values of hyperbolic elements of
$\Gamma$. This condition says that traces
may cluster, but only in a uniformly bounded
way. Bounded clustering is enough to force
exponential growth of the average
multiplicity.

\medskip\noindent
This is where the commutative algebra enters.
Given a Fuchsian group $\Gamma$, its invariant
trace field is
$$
K=
\mathbb{Q}
\bigl(\{\operatorname{tr}(\gamma^2):
\gamma\in\Gamma\}\bigr).
$$
The semiarithmetic condition says that $K$ is
totally real and that the traces are
``integers'' in $K$, meaning
$$
\operatorname{tr}(\gamma)\in\mathcal{O}_K.
$$
Here $\mathcal{O}_K$ is the ring of integers
of the number field $K$. This is exactly the
notion of integrality from commutative
algebra: an element is integral over a ring
if it satisfies a monic polynomial with
coefficients in that ring. Thus
$$
\mathcal{O}_K
=
\{x\in K:
x\text{ satisfies a monic polynomial in }
\mathbb{Z}[t]\}.
$$
Equivalently, $\mathcal{O}_K$ is the
integral closure of $\mathbb{Z}$ inside
$K$. This was the important conceptual point:
these are not ordinary integers, but they
are the elements of $K$ that behave like
integers relative to $\mathbb{Z}$.

\medskip\noindent
The reason this helps is that algebraic
integers in a fixed number field form a
lattice-like object. Since $K$ is totally
real, it has several real embeddings
$$
\sigma_i:K\hookrightarrow\mathbb{R}.
$$
So one trace has several real shadows
$\sigma_i(\operatorname{tr}\gamma)$. One of
these embeddings gives the actual Fuchsian
group acting on $\mathbb{H}^2$, while the
other embeddings provide arithmetic
constraints. In the arithmetic case there is
only one noncompact real factor. In the
semiarithmetic case there may be more than
one noncompact factor, and this number is
called the arithmetic dimension. The proof
uses the fact that the traces lie in
$\mathcal{O}_K$ and that their conjugates are
controlled. By a pigeonhole/lattice-counting
argument this gives bounded clustering of
trace values.

\medskip\noindent
A guiding example is the golden triangle
group
$$
\Gamma=\Delta(2,5,\infty).
$$
Its trace field is
$$
K=\mathbb{Q}(\sqrt{5}),
$$
with two real embeddings:
$$
\sigma_1(\sqrt{5})=\sqrt{5},
\qquad
\sigma_2(\sqrt{5})=-\sqrt{5}.
$$
The first embedding gives the actual
Fuchsian group. The second embedding gives
another real realization, which is not
discrete. This is typical of the
semiarithmetic situation: the group still has
arithmetic trace data, but it is not
arithmetic in the classical sense.

\medskip\noindent
The main story of the talk can therefore be
summarized as follows. Hyperbolic geometry
turns closed geodesics into conjugacy classes.
The length of a closed geodesic is determined
by a trace. If the traces live inside the
integer ring $\mathcal{O}_K$ of a number
field, commutative algebra gives arithmetic
control over how those traces can be
distributed. This control implies bounded
clustering, and bounded clustering forces
large average multiplicities in the length
spectrum. The new theorem says that this
mechanism works not only for arithmetic
surfaces, but also for certain
semiarithmetic surfaces.

\section{Widths, index, intersection
and isospectrality}
\label{section-widths-index-intersection}

\noindent
Jared Marx-Kuo, June 16, 2026,
IMPA, Differential Geometry Seminar.

\medskip\noindent
{\bf Abstract}
In this talk, I will discuss a series of
works on Gromov's $p$-widths,
$\{\omega_p\}$, on surfaces.
For ambient dimensions larger than $2$,
$\omega_p$ morally realizes the area of
an embedded minimal hypersurface of
index $p$. Historically, this was used
to prove the existence of infinitely
many minimal hypersurfaces in closed
Riemannian manifolds.
In ambient dimension $2$, the width
$\omega_p$ realizes the length of a
finite union of possibly immersed
closed geodesics:
$$
\omega_p
=
\sum_i m_i \ell(\gamma_i).
$$
Heuristically, one expects
$$
p
=
\sum_i \operatorname{Ind}(\gamma_i)
+
\#\{\text{vertices}\},
$$
where the vertices are the intersection
points of the union of geodesics.
Jointly with Lorenzo Sarnataro and
Douglas Stryker, we prove upper bounds
for the index and for the number of
vertices, making progress towards this
heuristic. Along the way, we prove a
generic regularity statement for
immersed geodesics.
If time allows, we will also discuss
the isospectral problem for the
$p$-widths, and explain how surfaces
provide a convenient setting to
investigate it.



\begin{definition}
\label{definition-minimal-hypersurface}
$Y^n$ is {\it minimal} if it's critical
for area, i.e. $\{Y_t\}$ with $Y_0=Y$,
$\frac{d}{dt}\Big|_{t=0}\operatorname{Area}
Y_t=0$, which is related to mean curvature,
$\leftrightarrow H_Y \equiv0$.
\end{definition}

\noindent
Applications:

\begin{itemize}
\item Inverse problems: how can we
relate the areas of the fibers of a
foliation with the ambient Riemannian
metric?

\item Physics: energy $\leftrightarrow$ area.
\end{itemize}

\noindent
The 1-dimensional version is geodesics.
We're interested in finding the curve
of minimal length in its homology class.

\medskip\noindent
{\bf Sweepouts.}
 $(M^{n+1},g)$ closed Riemannian manifold.
$$
\Phi:X^p \to Z_n(M)
$$
where we think of $Z_n(M)$ as the 
``space of hypersurfaces''.

\begin{example}
\label{example-1-sweepout}
$\Phi:[-1,1] \to Z_1(S^2)$ given by
$\Phi(t) = \{x = t\} \cap\{x^2+y^2+z^2=1\}$.
This is a foliation of $S^2$ by parallel
circles. For every point in the space,
there is a point in the domain
the interval, whose image is a fiber passing
through the given point.
\end{example}

\noindent
The full definition generalizes this
to $p$ points: if $\forall \{q_1,…,q_p\}$
there is $x \in X$ such that $\Phi(x)
\supseteq \{q_1,…,q_p\}$.

\begin{example}
\label{example-rp2-sweepout}
$\Phi:\mathbb{R}P^2 \to Z_1(S^2)$,
$\Phi[a:b:c] = \{ax+by+cz=0\}\cap S^2$.
is a 2-sweepout.
\end{example}

\noindent
Now consider
$$
P_p = \{\Phi:X^p \to Z_n(M):
\Phi \operatorname{sweepout}\}
$$
\begin{definition}
\label{definition-gromov-p-width}
Let $p$ be a positive integer.
The {\it Gromov $p$-width} is
$$
\omega_p(M,g) = \text{inf}_{\Phi \in P_p}
\text{sup}_{x \in \operatorname{Dom}(\Phi)}
\operatorname{Area}(\Phi(x)).
$$
\end{definition}

\noindent
Which Gromov explained is some sort
of ``non-linear spectrum of area
functional''.
It is positive and, $\omega_p \leq W_{p+1}$.

\begin{theorem}
\label{theorem-weyl-sweepouts}
$$
\lim_{p \to \infty}\lambda_p
p^{\frac{-2}{n+1}}
= \frac{\alpha_{n+1}}{\text{Vol}(M,g)}.
$$
\end{theorem}

\begin{theorem}[GGLMN]
\label{theorem-gglm}
$$
\lim_{p \to \infty}\omega_p
p^{\frac{-1}{n+1}}
=B_{n+1}\text{Vol}(M,g)^{\frac{n}{n+1}.}
$$
$\omega_p$ give existence of minimal surfaces.

\end{theorem}

\noindent
The idea is to see how this number
grows as $p$ varies.
The upshot is that this connects
with minimal surfaces!

\begin{theorem}[M-Neves-Z]
\label{theorem-mnz}
$(M^{n+1},g)$, $3 \leq  n+1 \leq  ?$,
$g$ ``generic'', then
$$
\omega_p(M,g)
=\operatorname{Area}(\Sigma_p),
$$
$\Sigma_p$ minimal embedded with multiplicity
1.
\end{theorem}

\begin{lemma}[Corollary, Yau's conjecture]
\label{lemma-yau-conjecture-minimal-surfaces}
Every closed $(M^{n+1},g)$,
$3 \leq  n+1 \leq ?$
has infinitely minimal surfaces.
\end{lemma}

\noindent
Other results:
\begin{itemize}
\item Existence of many CMCs (D,M-Z).
\item Existence of many PMCs (G-MK).
\end{itemize}

\noindent
We look for objects realising the $p$-widths.

\noindent
Questions:
\begin{enumerate}
\item What about $1=2$?
\item Can we compute $\omega_p$ for nice
$(M,g)$?
\item When does $\{\omega_p(M,g)\}_{p=1}^\infty$
determine $(M,g)$?
Can we hear the shape of a drum?
\end{enumerate}

\begin{theorem}[C-M]
\label{theorem-cm1}
$(M^2,g)$
closed smooth. $\forall  p$ there exists
$\{\gamma_p\}_{p=1}^{N_p}$
closed immersed geodesics
$\{m_p\}_{p=1}^{N_p}\subset \mathbb{Z}^+$
such that
$$
\omega_p(M^2,g)
=\sum_{i=1}^{N_p}m_{i,p}\ell(\gamma_{i,p}.
$$
\end{theorem}

\begin{theorem}[C-M]
\label{theorem-cm2}
$$
\omega_p(S^2,g_{\operatorname{round}})
=2\pi \operatorname{floor}(\sqrt{p}).
$$
\end{theorem}

\noindent
Contrast:
\begin{theorem}[Ambrosio-M-N]
\label{theorem-amn}
$(\mathbb{R}P^2,g_{\operatorname{round}})$
is $\omega_p$-spectrally rigid, that is,
if $(M^{n+1},g)$ is such that
$\omega_p(M^{n+1},g)=\omega_p(\mathbb{R}P^2,g_{\operatorname{round}})$
for all $p$ then
$(M^{n+1},g)$ is isometric to
$(\mathbb{R}P^2,g_{\operatorname{round}})$.
\end{theorem}

\medskip\noindent
{\bf Zoll metrics.}

 \begin{definition}
\label{definition-zoll-metric}
A {\it Zoll metric} is such
that all geodesics are simple 
(they do not autointersect) and of the
same length.
\end{definition}

\noindent
There is a Zoll metric on $S^2$ which is
not the round metric, found by Zoll
solving a very particular ODE.

\begin{theorem}[G-F]
\label{theorem-gf-moduli}
For a fixed positive number $\ell_0$
there is an infinite-dimensional moduli
space of Zoll metrics on $S^2$, $Z$
with
$$
T_{g_{\operatorname{round}}}Z
\cong
\{f:S^2 \to \mathbb{R}:
f(-x) = -f(x)\}.
$$
\end{theorem}

\medskip\noindent
{\bf Index and intersection on surfaces.}
Recall that for $3 \leq  n+1 \leq  ?$,
$$
\omega_p = \operatorname{Area}(\Sigma_p)
$$
and
$$
p = \operatorname{Ind}(\Sigma_p).
$$

Note that the latter condition about
the index cannot be true for $n=2$.

The idea is that points of intersection
amog $\bigcup_{i=1}^{N_p}\gamma_{i,p}$
constribute to ``generalized index'' $=p$.

Conjecture: if $m_{i,p}=1$,
$$
p=\sum_{i=1}^{N_p}\operatorname{Ind}(\gamma_{i,p}+
\alpha
$$
where $\alpha$ is the number of vertices
(i.e. intersections)
of $\bigcup_{i=1}^{N_p}\gamma_{i,p}$.

\begin{theorem}[MK-S-S]
\label{theorem-mkss}
If $\omega_p =
\sum_{i=1}^{N_p}m_{i,p}(\gamma_{i,p})$
with $\gamma_{i,p}$ immersed geodesics,
then
$$
\sum_{i=1}^{N_p}\operatorname{Ind}(\gamma_{i,p})
\leq p,
$$
$$
\alpha \leq  p,
$$
where $\alpha$ is as above.
\end{theorem}

\begin{theorem}
\label{theorem-geodesics-intersect}
For generic $g$ on $M^2$,
any finite union of geodesics has
intersections
whcih look like 2 lines at a point.
\end{theorem}

\section{Singular Riemannian Foliations,
Variational Problems and Principles of
Symmetric Criticalities}
\label{section-singular-riemannian-foliations}

\medskip\noindent
{\bf Conference information.}

\href{%
https://sites.google.com/view/%
60thanniversary-dmat-puc-rio/Home
}{International Conference Celebrating
the 60th Anniversary of the Department
of Mathematics at PUC-Rio}

August 12--14, 2026,
RDC Auditorium, PUC-Rio,
Rio de Janeiro.

\medskip\noindent
{\bf Speaker.}

Marcos Alexandrino, USP.

\medskip\noindent
{\bf Date and time.}

August 13, 2026, 13:30--14:30.

\medskip\noindent
{\bf Title.}

Singular Riemannian Foliations,
Variational Problems and Principles of
Symmetric Criticalities.

\medskip\noindent
{\bf Abstract.}

A singular foliation on a complete
Riemannian manifold $M$ is called a
singular Riemannian foliation if its
leaves are locally equidistant. An example
is the partition of $M$ into the orbits of
an isometric action. The talk investigates
problems in the calculus of variations on
compact Riemannian manifolds equipped with
singular Riemannian foliations having
special properties. The examples under
consideration include isoparametric
foliations, singular Riemannian foliations
on Euclidean fiber bundles, and partitions
into the orbits of Lie groups acting by
isometries. This is joint work with
Leonardo F. Cavenaghi, Diego Corro, and
Marcelo K. Inagaki. The talk is intended
for a general audience interested in
differential geometry and analysis.


\medskip\noindent
{\bf Main goals.}

\begin{enumerate}
\item[(I)] Present the new class of
singular foliations adapted to
variational problems (AVP). This class
should include homogeneous foliations,
as well as nonhomogeneous examples such
as isoparametric foliations and SRF on
Euclidean fiber bundles.

\item[(II)] Define the new concept of
$\mathcal F$-symmetric criticality in
order to include interesting operators
$J:W^{1,p}(M)\to\mathbb R$, such as
$$
\begin{aligned}
J_\lambda(f)
&=
\mathcal M(\|f\|^p)
\int_M
\mathcal L(
|\nabla f|^2,f,x)\,\mathrm{vol}
\\
&\quad
-\frac{\lambda}{a}
\left(
\int_M F(f,x)\,\mathrm{vol}
\right)^{r+1}.
\end{aligned}
$$
We also aim to establish
$\mathcal F$-principles of symmetric
criticality.

\item[(III)] Apply the
$\mathcal F$-principles to find
infinitely many weak solutions
$\{u_i\}_{i=1}^{\infty}$ of the
variational problem
$$
dJ_\lambda(u_i)=0.
$$
\end{enumerate}

\end{document}
